\documentclass[11pt]{article}
\usepackage[letterpaper,margin=1.05in]{geometry}
\usepackage[T1]{fontenc}
\usepackage{lmodern,amsmath,amssymb,amsthm,microtype,booktabs,array,tabularx}
\usepackage[round,authoryear]{natbib}
\usepackage[colorlinks=true,linkcolor=blue,citecolor=blue,urlcolor=blue]{hyperref}
\newtheorem{theorem}{Theorem}
\newtheorem{proposition}[theorem]{Proposition}
\newtheorem{lemma}[theorem]{Lemma}
\newtheorem{corollary}[theorem]{Corollary}
\theoremstyle{definition}
\newtheorem{definition}[theorem]{Definition}
\newtheorem{remark}[theorem]{Remark}
\newtheorem{example}[theorem]{Example}
\newcommand{\GH}{\mathsf{GH}}
\newcommand{\Mink}{\mathbb M}
\newcommand{\U}{\mathcal U}
\newcommand{\Ifut}{I^{+}}
\newcommand{\Ichr}{I^{-}}
\newcommand{\Jp}{J^{+}}
\newcommand{\Jm}{J^{-}}
\newcommand{\B}{\mathcal B}
\newcommand{\Is}{I^-_{\mathrm s}}
\newcommand{\Isp}{I^+_{\mathrm s}}
\newcommand{\Det}{\mathrm{Det}}
\newcommand{\partheading}[1]{\bigskip\begin{center}\large\bfseries #1\end{center}\addcontentsline{toc}{section}{#1}}
\newcolumntype{Y}{>{\raggedright\arraybackslash}X}
\title{What the Past Can Fix:\\[2pt]
\large becoming, determination and existence in general relativity}
\author{Andrew Impellitteri}
\date{31 October 2026 (working draft)}
\begin{document}
\maketitle

\begin{abstract}
Four questions about the future are often run together: whether it has become,
whether it is fixed, whether it exists, and whether claims about it are already
true. We separate them, with theorems where the question is mathematical and with countermodels to stated inferences where it is not. Under a locality premise, what has become
for an event in any globally hyperbolic spacetime lies between its chronological
and causal past, unless everything has become for everything, and a transitive local relation of becoming defined from the metric alone admits no shared present. What is fixed
from an event is a different set. It contains what the causal structure
guarantees, which we characterize by slices and compute in the standard models,
and a particular law can fix more. In the closed dust model a massless field is
fixed on part of an event's far future while its near future stays open, and every
massive field is fixed exactly on the past before maximum expansion and
everywhere from then on. In a smooth complete class of Einstein--scalar
spacetimes, conditional on named external estimates, the whole pointed causal
past and the total incoming mass leave an open ball of future central readings
free, while agreement on everything outside the causal future fixes the entire
development. So whether the future is open depends on which base is held fixed and on the law. Given a base and a law, the answer is the same whatever one says about whether the future exists; if the base is taken to be what exists from an event, the answer varies with the ontology. Fixing also has no arrow of its own: for time-reversal invariant laws every such result has a mirror image, so the slogan that the past is fixed and the future open needs another source. We then ask what would be needed to reach existence. Stated bridges about constitution and persistence supply lower bounds on what is concrete from an event. A restricted ontology needs an upper bound, an
unrestricted one a stronger lower bound, and neither comes from the results
above. We sort the stated candidates, and record which inferences fail.
No temporal ontology is selected.
\end{abstract}

\tableofcontents

\section{Introduction}\label{sec:intro}

Is the future real? Is it already fixed? Is what we say about it already true?
Has it become? In ordinary thought these are nearly one question. An open future
is one that does not yet exist, whose outcome nothing settles, about which there
is no fact yet. Philosophers of time have long known that the pieces can be
pulled apart in principle. A growing block in Broad's sense \citep{Broad1923}, with deterministic laws added, has an unreal future that is settled. A later-grounded bivalent theory has
truths about what does not yet exist \citep{CR2018}. What has been missing is an
account, in general relativity, of how far apart the pieces actually lie, and of
exactly what fixes what.

This paper gives that account. It has one thesis:
\begin{quote}
What has become for an event, what is fixed from it, what exists from it and what
is true from it are four different relations. General relativity constrains the
first sharply under a locality premise. It makes the second a mathematical question, once a law and a comparison are named, and we answer it exactly in the models treated. Nothing here determines the third or the fourth, and the first two do not determine them.
\end{quote}

The thesis is defended by results of four kinds. We state them here with their
conditions.

\paragraph{Becoming (Part I).} Stein's theorem, with its contextual conditions, says that in Minkowski spacetime a transitive, symmetry-invariant relation of becoming that contains the causal past and leaves something unsettled is the causal past \citep{Stein1968,Stein1991}. Clifton and Hogarth give the form whose alternatives are the chronological past, the causal past and the universal relation \citep{CH1995}. A curved spacetime can have no
symmetries, and the single-spacetime statement is then false. Replacing symmetry
by the local covariance of quantum field theory restores the result: either
everything has become for everything in every globally hyperbolic spacetime, or in
every such spacetime becoming lies between the chronological and causal past
(Theorem~\ref{thm:main}). The residual freedom is the light cone, and it can depend
on curvature. For transitive relations that contain the chronological past and are defined from the metric alone, locality, a shared present and non-trivial becoming cannot all be had (Corollary~\ref{cor:nopresent}). Without transitivity they can: the bow tie of Section~\ref{sec:rigidity} has all three.
A rival with a preferred clock changes the category and is not excluded
(Section~\ref{sec:becoming-rival}).

\paragraph{Determination on fixed backgrounds (Part II).} What the past fixes is a
different set. Theorem~\ref{thm:char} characterizes what the causal structure
guarantees, by a test on slices. Theorem~\ref{thm:map} then lays out what can
happen. The fixed set can equal the causal past, exceed what is guaranteed on an
isolated part of the far future, or be everything; the same metric gives different
answers under different laws; and an exact verdict can flip while the hidden fields change by $O(m^2)$ in a stated norm (equation~\eqref{eq:hidden}). In one closed universe a massless field is fixed at the end of an
event's worldline before the middle (Theorem~\ref{thm:diamonds}), and every massive
field has no such islands (Theorem~\ref{thm:massive}).

\paragraph{Determination with gravity (Part III).} Backreaction does not close the
gap. In a smooth, complete, inextendible class of Einstein--scalar spacetimes, the
entire pointed causal past of an event and the total incoming mass leave an open
ball of future central scalar readings free (Theorem~\ref{thm:same-past}), yet
agreement on the open region outside the event's causal future fixes the whole
development (Theorem~\ref{thm:rigidity}). Both statements concern the same class. The
first is conditional on external estimates of Luk, Oh and Yang. Its moral is already
present for the linear wave equation; what coupling adds is that the moral survives
gravity, a fixed mass and the axis. Which base is compared therefore matters more than
any general slogan about openness.

\paragraph{Existence and truth (Part IV).} None of this reaches existence. We show
what would: a stated principle bounding what is concrete from an event. Stated bridges from the physics give lower bounds. A restricted ontology needs an upper bound, and an unrestricted
one needs a stronger lower bound (Section~\ref{sec:bounds}). We sort the stated candidates. The useful result is a diagnostic for these candidates: an index on ``concrete from $b$'' is either defined, and then trivial and shared, or primitive, and then a thesis the event realist rejects. Three candidates need support of a different kind. Truth policies are independent again (Section~\ref{sec:truth}). Fixing has no arrow of its own, so the past as a base needs a separate account (Section~\ref{sec:direction}). Section~\ref{sec:revise} records which inferences fail.

\paragraph{What is not claimed.} No temporal ontology is selected. The results are
conditional as stated, and Appendix~\ref{app:status} gives the status of each. Parts
of this material overlap prior work, which is credited where it is known; a claim of
priority would need an audit this paper has not completed.

\section{Setting}\label{sec:setting}

Spacetimes are smooth, connected, four-dimensional, oriented and time-oriented.
Unless stated otherwise they are globally hyperbolic. We follow Clifton and Hogarth
in taking $\Ichr,\Jm$ to be reflexive \citep[p.~360]{CH1995}. We write $\Is(p)$ for the
strict chronological past, so $\Ichr=\Is\cup\{p\}$, $x\ll y$ for $x\in\Is(y)$, and
$x\le y$ for $x\in\Jm(y)$. Standard causal facts are cited from Minguzzi and Sánchez
\citep{MS2008}, abbreviated MS, whose $I^\pm$ are our strict sets. A \emph{Cauchy time
function} is a continuous surjection $t:M\to(\alpha,\beta)$ strictly increasing along
future causal curves, with Cauchy level sets $S_c=t^{-1}(c)$ \citep{BS2005}.

Fix an event $b$, which we call the standpoint. We attach to it the following.

\begin{center}
\begin{tabularx}{\textwidth}{@{}lYl@{}}
\toprule
Object & Meaning & Decided by\\
\midrule
$B_M(b)$ & what has become for $b$ & a locality premise (Part I)\\
$\Det(b)$ & what the causal structure guarantees from $\Is(b)$, fixed by every law with the domain-of-dependence property & causal structure (Part II)\\
$F_{\mathcal L}(S)$ & $x\in F_{\mathcal L}(S)$ if any two solutions agreeing on $S$ agree at $x$; $F_{\mathcal L}(b):=F_{\mathcal L}(\Is(b))$ & law, class and base (Parts II, III)\\
$E_b$ & what is concrete from $b$ & a bridge premise (Part IV)\\
truth at $b$ & which claims about later events hold from $b$ & a semantic policy (Part IV)\\
\bottomrule
\end{tabularx}
\end{center}

These are not five names for one thing, and the paper is about the gaps between them.
Three points of method follow from keeping them apart.

\emph{First, name the comparison.} Agreement of accessible records, agreement of an
intrinsic causal past, and agreement of an entire proposed realm are three different
contracts. A failure under one is not a failure under another. Here we mostly compare exact agreement of classical fields on a named region.

\emph{Second, separate physical premises, definitions and bridges.} A definition that
already assumes an ontology cannot establish it. We mark bridge premises as such.

\emph{Third, keep models in their domains.} A finite example or a closed dust universe
shows what is possible under stated hypotheses. Neither is a claim about the actual
cosmology.

\partheading{Part I. Becoming}

\section{Becoming under locality}\label{sec:becoming}

\begin{definition}
$\GH$ is the category whose objects are globally hyperbolic spacetimes and
whose morphisms $\psi:M\to M'$ are smooth isometric embeddings preserving
orientation and time orientation, with $\psi(M)$ open and causally convex in
$M'$. This is the category $\mathsf{Loc}$ of locally covariant quantum field
theory \citep[\S2]{BFV2003}, \citep[\S2.2]{FV2012}.
\end{definition}

\begin{definition}
A \emph{becoming theory} assigns to each $M\in\GH$ a relation
$B_M\subseteq M\times M$ (read $x\,B_M\,b$ as ``$x$ has become for $b$'') such
that, for every $M$:
\begin{enumerate}
\item[(LC)] for every morphism $\psi:M\to M'$, $x\,B_M\,b\iff
\psi(x)\,B_{M'}\,\psi(b)$;
\item[(T)] $B_M$ is transitive;
\item[(CB)] $\Ichr_M\subseteq B_M$.
\end{enumerate}
\end{definition}

(CB) is Clifton and Hogarth's chronological becoming; they derive it from
``worldline becoming'' and ``worldline transitivity'' \citep[Lemma,
p.~364]{CH1995}, and the derivation carries over verbatim to curved
spacetime. (LC) contains invariance under every orientation- and
time-orientation-preserving isometry of each $M$, which on Minkowski spacetime
$\Mink$ is the proper orthochronous Poincaré group. Its further content is
\emph{locality}: if $N\subseteq M$ is open, causally convex and contains $x$
and $b$, then whether $x$ has become for $b$ is settled by $N$ alone.

\begin{remark}[What (LC) contains]\label{rem:lc}
Three premises should be kept apart. \emph{Naturality}: one rule for all
objects, preserved by every morphism. \emph{Reflection}: the converse
direction of (LC), by which every causally convex region containing $x$ and
$b$ settles the verdict; the proof of Theorem~\ref{thm:main} uses it along $N\hookrightarrow M$ and $N''\hookrightarrow M'$ in part (a), and along $\psi|_N:N\to M'$ and $N''\hookrightarrow\Mink$ in part (b), and functoriality in quantum
field theory supplies no such direction. \emph{Off-shell universality}:
Remark~\ref{rem:offshell}. ``Local'' here is geometric: independence of what
lies around a region, the sense of \citet[\S1]{BFV2003}. It is not
no-signalling or microcausality, and neither of those implies it. The case for
(LC) is that a relation defined from geometry alone should not depend on what
lies outside a causally convex region containing the related events. That is
a premise about becoming, not a consequence of relativity.
\end{remark}

\begin{remark}[The off-shell premise]\label{rem:offshell}
$\GH$ contains every globally hyperbolic metric, not only solutions of
Einstein's equations; the category is defined kinematically in
\citep{BFV2003,FV2012}. Theorem~\ref{thm:main} uses this: its proof builds
auxiliary spacetimes that need not satisfy any field equation with reasonable
matter. This is natural for a kinematic notion of becoming, defined (as Stein
and Clifton--Hogarth define theirs) from geometry alone, but it is a premise.
Section~\ref{sec:energy} examines what happens without it.
\end{remark}


\subsection{Why one spacetime is not enough}\label{sec:single}

\begin{proposition}\label{prop:single}
Let $M$ be globally hyperbolic and $x,b\in M$ spacelike-separated. Then
$R=\Jm\cup(\Jm(x)\times\Jp(b))$, where $(e,z)$ in the second set means $e$
has become for $z$, is reflexive and transitive, contains $\Ichr$, relates $x$
to $b$, and is neither $\Ichr$, $\Jm$ nor $\U$. If $M$ has no nontrivial
isometries, $R$ is invariant under all of them.
\end{proposition}

\begin{proof}
Write $P=\Jm(x)\times\Jp(b)$. Compositions $\Jm\circ\Jm$, $\Jm\circ P$ and
$P\circ\Jm$ land in $\Jm$ or $P$ because $\le$ is transitive. A composition
$P\circ P$ would need some $q$ with $b\le q\le x$, which is impossible. The
pair $(b,x)$ is not in $R$, so $R\neq\U$.
\end{proof}

So for a single generic spacetime, Stein-type uniqueness fails completely. The
next example shows that
uniformity across spacetimes is still not enough without locality.

\begin{proposition}\label{prop:nonlocal}
Let $\tau_M(p)$ be the supremum of the lengths of past-directed causal curves
from $p$ (the cosmological time of Andersson, Galloway and Howard
\citep{AGH1998}). Put $B_M=\{(x,b):\tau_M(x)\le\tau_M(b)\}$ if $\tau_M$ is
finite on $M$, and $B_M=\Jm_M$ otherwise. Then $B$ is invariant under all
isometries, each $B_M$ is
reflexive and transitive and satisfies (CB), and $B$ relates spacelike pairs
in every spacetime where $\tau$ is finite (e.g.\ big-bang
Friedmann--Lemaître--Robertson--Walker models). It is not a becoming theory:
it violates (LC).
\end{proposition}

\begin{proof}
Isometries preserve lengths of causal curves. If $x\ll b$, concatenating
past curves from $x$ with a timelike curve from $x$ to $b$ gives
$\tau(b)\ge\tau(x)+d(x,b)>\tau(x)$, so (CB) holds; transitivity is that of
$\le$ on the reals. Failure of (LC): the Milne wedge $W=\Ifut_\Mink(0)$ is open,
causally convex and globally hyperbolic, and its inclusion into $\Mink$ is a
morphism. On $W$, $\tau_W(p)$ is the proper time from the origin, finite, so
$B_W$ relates distinct points of the same hyperboloid $\{\tau_W=c\}$, which are
spacelike-separated; on $\Mink$, $\tau=\infty$ and $B_\Mink=\Jm$ does not relate
them.
\end{proof}


\subsection{The classification}\label{sec:results}

\begin{theorem}[Clifton--Hogarth]\label{thm:flat}
For every becoming theory, $B_\Mink\in\{\Ichr,\Jm,\U\}$.
\end{theorem}

\begin{proof}
Proper orthochronous Poincaré maps are automorphisms of $\Mink$ in $\GH$, so
$B_\Mink$ is invariant under them by (LC). Clifton and Hogarth's Theorem~1
\citep[p.~365]{CH1995} gives the conclusion from this invariance,
transitivity and chronological becoming. Its proof (pp.~365--367) uses only
point-preserving spatial rotations, boosts and translations, and they note
that it does not use spatial reflections (pp.~385--386). Their stated
hypothesis is invariance under all time-orientation-preserving isometries,
reflections included; using only the proper orthochronous subgroup is a
weakening whose adequacy rests on their proof, not their statement.
\end{proof}

\begin{theorem}[No becoming outside the causal past]\label{thm:main}
Let $B$ be a becoming theory.
\begin{enumerate}
\item[(a)] If $B_\Mink\neq\U$, then $B_M\subseteq\Jm_M$ for every $M\in\GH$.
\item[(b)] If $B_\Mink=\U$, then $B_M=\U_M$ for every $M\in\GH$.
\end{enumerate}
Hence either $B$ is universal everywhere, or $\Ichr_M\subseteq
B_M\subseteq\Jm_M$ for every $M$.
\end{theorem}

Call $B$ \emph{universal} if $B_M=\U_M$ for every $M$. By
Theorem~\ref{thm:main} this holds if and only if $B_\Mink=\U$, and if $B_M=\U_M$
for one $M$ then it holds for all; otherwise $B$ is \emph{non-universal}.

\begin{corollary}[No shared present]\label{cor:nopresent}
Call $x$ and $b$ \emph{co-present} in $M$ if $x\,B_M\,b$ and $b\,B_M\,x$. For
every non-universal becoming theory, co-presence is the identity relation in
every $M\in\GH$. Equivalently, no becoming theory is local, non-trivial and
admits two distinct co-present events in some globally hyperbolic spacetime.
\end{corollary}

\begin{proof}
By Theorem~\ref{thm:main}(a), $x\le b$ and $b\le x$, and globally hyperbolic
spacetimes are causal, so $x=b$.
\end{proof}

\begin{corollary}[Worldline-dependent becoming]\label{cor:worldline}
Suppose that for each $M\in\GH$, each inextendible timelike curve $\lambda$ in
$M$ and each $b\in\lambda$ a set $B_{M,\lambda}(b)$ is given, satisfying
Clifton--Hogarth's worldline becoming ($\Ichr(b)\cap\lambda\subseteq
B_{M,\lambda}(b)$) and worldline transitivity ($a\in B_{M,\lambda}(b)$ and
$b\in B_{M,\lambda'}(c)$ imply $a\in B_{M,\lambda'}(c)$) \citep[pp.~362--363]{CH1995},
and local covariance in the form: for a morphism $\psi:N\to M$, an inextendible
timelike $\lambda_N$ in $N$ and any inextendible timelike $\lambda$ in $M$
extending $\psi\circ\lambda_N$, $a\in B_{N,\lambda_N}(b)\iff\psi(a)\in
B_{M,\lambda}(\psi(b))$. Then $B_{M,\lambda}(b)$ does not depend on $\lambda$, and
the resulting relation is a becoming theory. Hence Theorem~\ref{thm:main}
applies to worldline-dependent becoming.
\end{corollary}

Clifton and Hogarth anticipated that extending their Theorem~2 would be
harder still (``much less to extend Theorem 2'', p.~385). Under local
covariance it is not: as they themselves note (p.~363), worldline transitivity
makes observers whose worldlines cross agree, and the remaining work is done by
Theorem~\ref{thm:main}.

Write $E^-=\Jm\setminus\Ichr$ for the horismos (lightlike-only relation).

\begin{theorem}[Classification of the residual freedom]\label{thm:class}
Let $B$ be a non-universal becoming theory and $N_M=B_M\cap E^-_M$. Then $N$
is locally covariant and satisfies
\begin{enumerate}
\item[(NC)] if $(x,y),(y,b)\in N_M$ and $(x,b)\in E^-_M$, then $(x,b)\in N_M$.
\end{enumerate}
Conversely, if $N\subseteq E^-$ is locally covariant and satisfies (NC), then
$B=\Ichr\cup N$ is a becoming theory. On $\Mink$, $N$ is empty or all of $E^-$.
\end{theorem}

\begin{theorem}[Non-uniqueness]\label{thm:weyl}
Let $N_M(x,b)$ hold iff $x\in E^-(b)$, $x\neq b$, and some lightlike geodesic
segment from $x$ to $b$ contains a point at which the Weyl tensor is nonzero.
Then $\Ichr\cup N$ is a becoming theory equal to $\Ichr$ on $\Mink$, equal to
$\Jm$ on any spacetime with nowhere-vanishing Weyl tensor (such as the
Schwarzschild exterior), and strictly between $\Ichr$ and $\Jm$ on spacetimes
that are flat on a nonempty causally convex open set and have nonzero Weyl
tensor somewhere. It depends only on the conformal structure (not merely the
curvature: conformally flat spacetimes such as FLRW give $\Ichr$).
\end{theorem}

\begin{proposition}[Regularity]\label{prop:reg}
Let $B$ be non-universal. If every $B_M$ is closed in $M\times M$, then
$B=\Jm$. If every $B_M\setminus\Delta_M$ is open, then $B=\Ichr$.
\end{proposition}


Proofs are in Appendix~\ref{sec:proofs}.

\subsection{Discrete and extreme cases}\label{sec:variants}

\subsection{Causal sets}\label{sec:causets}

Let $\mathsf{CS}$ have locally finite posets as objects and order-embeddings
with order-convex image as morphisms. A becoming theory on $\mathsf{CS}$ is
defined as before, with (CB) read as $x\le b\Rightarrow x\,B\,b$.

\begin{proposition}\label{prop:causets}
Every becoming theory on $\mathsf{CS}$ is either the causal order in every
causal set or the universal relation in every causal set.
\end{proposition}

\begin{proof}
If $x,b$ are incomparable, $\{x,b\}$ is convex, so both embeddings of the
two-element antichain are morphisms: the verdict on incomparable pairs is a
single symmetric bit $\beta$, the same in every causal set. If $x>b$, the
interval $I=[b,x]$ is finite and convex, so $x\,B\,b$ depends only on $I$. If
$\beta=1$, in $I\sqcup\{z\}$ with $z$ incomparable to all, $x\,B\,z\,B\,b$; so
every such verdict is positive and $B$ is universal. If $\beta=0$ and
$x\,B\,b$ for some $x>b$, adjoin $e$ with $e<x$ only: $I$ remains convex, so
$x\,B\,b$ there, $e\,B\,x$ by (CB), and $e\,B\,b$ with $e$ incomparable to $b$,
contradicting $\beta=0$.
\end{proof}

Wüthrich and Callender state that ``The analogue of this theorem, however, is
straightforwardly false in CST'' \citep[\S3]{WC2017}. Their counterexample, on
the diamond causal set $d<b,c<a$, adds the incomparable pair $\{b,c\}$ to the
causal order and is invariant under that causal set's automorphisms. It is not
the restriction of any becoming theory in the present sense: relating $b$ to
$c$ fixes $\beta=1$, and Proposition~\ref{prop:causets} then relates $b$ to $d$,
which their relation excludes. Their own footnote~5 observes that the
counterexample ``could be a small proper part of a much larger causet''; (LC) is
exactly the requirement that the relation on such a part agree with the
relation in the larger causal set. The diagnosis matches the continuum case:
invariance within one structure is too weak, and locality across structures
restores Stein's theorem, here with no residual freedom because a causal set
has no light-cone boundary. Their claim is true on its own terms; the present
result uses a stronger, cross-structure hypothesis. Note also that
$\mathsf{CS}$ includes posets that are not manifoldlike (antichains, disjoint
unions), and the proof uses them: this is the causal-set counterpart of the
off-shell premise of Remark~\ref{rem:offshell}.

\subsection{Gödel universes are isolated}\label{sec:godel}

\begin{proposition}\label{prop:godel}
Extend $\GH$ to all oriented, time-oriented spacetimes with the same
morphisms. If $M$ is totally vicious ($p\in\Ichr_{\mathrm s}(p)$ for all $p$),
then every morphism into $M$ is onto, every theory satisfying (CB) has
$B_M=\U_M$, and there is no morphism between $M$ and any causal spacetime, in
either direction.
\end{proposition}

\begin{proof}
In a totally vicious spacetime $\Ichr(p)=\Ifut(p)=M$ for every $p$ (MS \S3.1).
Let $V$ be nonempty, open and causally convex, $p\in V$ and $q\in M$. There are
timelike curves $p\to q$ and $q\to p$; their concatenation is a causal curve
with both endpoints in $V$, so it lies in $V$, and $q\in V$. Thus $V=M$
(compare MS Remark~2.12). (CB) gives $B_M\supseteq\Ichr_M=M\times M$. A morphism
$N\to M$ from a causal $N$ would be onto, hence an isometry, which is
impossible since $M$ has closed timelike curves; a morphism $M\to N$ would map
a closed timelike curve of $M$ to one of $N$.
\end{proof}

Gödel's universe is totally vicious, so universal becoming there is forced;
this is Gödel's first step, in Stein's terms. But local covariance transmits
constraints only along morphisms, and by Proposition~\ref{prop:godel} no
morphism connects a totally vicious spacetime with a globally hyperbolic one. Gödel's modal step, from what holds
in such a universe to what holds in ours \citep[pp.~561--562]{Godel1949}, as
quoted in \citep[\S6]{Monton2006}, requires that becoming be the same
feature across nomologically possible universes. Local covariance is one
precise way of requiring a rule to be the same across universes, and it
transmits nothing out of totally vicious ones. A version of Gödel's step that does go through
would need a different cross-universe link. Proposition~\ref{prop:nonlocal} shows what such links
look like: natural but non-local. Compare the conceptual responses of
\citet{Savitt2005}, \citet{Dieks2006} and \citet{Aames2022}, which locate
becoming locally along worldlines.


\subsection{Energy conditions}\label{sec:energy}

Let $\GH_{\mathrm{DEC}}$ be the full subcategory of spacetimes whose Einstein
tensor satisfies the dominant energy condition. It is closed under restriction
to causally convex open subsets and contains $\Mink$.

\begin{proposition}[The auxiliary spacetime is unavailable]\label{prop:dec}
If $M\in\GH_{\mathrm{DEC}}$ is flat on a neighbourhood of some totally geodesic
Cauchy surface, then $M$ is flat. In particular the spacetime $M'$ of Appendix~\ref{sec:proofs} lies in $\GH_{\mathrm{DEC}}$ only if
$g$ is flat on $T_{2r}$.
\end{proposition}

\begin{proof}
On the Cauchy surface $S$ the metric is flat and the Einstein tensor vanishes. By the conservation theorem
for matter obeying the dominant energy condition \citep[\S4.3]{HE1973}, the
Einstein tensor vanishes on its domain of dependence, i.e.\ everywhere. In the
form we use: let $t$ be a Cauchy temporal function with $S=t^{-1}(0)$
\citep[Thm.~1.2]{BS2006} and $p\in J^+(S)$. The set $J^-(p)\cap J^+(S)$ is compact, and
it is the union of the segments of integral curves of $\nabla t$ that run from
$S$ to $\partial J^-(p)$ (a future-directed curve that leaves the past set
$J^-(p)$ does not return). Minguzzi's form of the theorem \citep[Thm.~1]{Minguzzi2015},
which needs no time function, then gives $G=0$ at $p$ from $G=0$ on $S$;
$J^-(S)$ is the time reverse. For
$M'$, any slice $\{t=c\}$ with $c<-3r$ is totally geodesic and has a flat
neighbourhood.
A vacuum globally hyperbolic spacetime with flat data $(\Sigma,h,0)$ on a
Cauchy surface is a development of those data and so embeds isometrically in
their maximal development \citep[Thm.~3]{CBG1969}. That maximal development is
the flat static spacetime $\mathbb R\times\Sigma$, $-dt^2+h$ (not Minkowski in
general; e.g.\ $\Sigma=T^3$): it is a development, and if a development
properly extended it, a point $p$ outside its image would lie on a timelike
geodesic from $\Sigma$ whose image under the extension coincides with the
corresponding complete geodesic of $\mathbb R\times\Sigma$, so $p$ would lie in
the image. Hence the spacetime is flat.
\end{proof}

The same reasoning constrains any repair that keeps the original spacetime and
inserts small flat pockets near the points $w$, $v$ and the detour of the proof of Theorem~\ref{thm:main}: in a vacuum region
whose causal past on some Cauchy surface is vacuum and unmodified, the
conservation theorem forces the modified spacetime to be vacuum there with the
same data, so no such pocket can be inserted there. Pockets must be fed by matter brought
in from the past and routed around $D(\Sigma)$ afterwards. The proof of
Theorem~\ref{thm:main} goes through verbatim in $\GH_{\mathrm{DEC}}$ given
the following sufficient condition, which we leave open.

\begin{quote}
\textbf{Pocket problem.} Given $M\in\GH_{\mathrm{DEC}}$ and $\Sigma$ as in Lemma~\ref{lem:D},
find $M''\in\GH_{\mathrm{DEC}}$ containing an isometric copy of a neighbourhood
of $D_M(\Sigma)$ whose image is a domain of dependence, and exactly flat
causally convex regions containing points $w\ll x$ and $v\gg b$ joined by the
flat domain of dependence of a slab.
\end{quote}

The constraint is an inequality, not a field equation, so constructions can be
kinematic, and explicit models with flat regions chronologically adjacent to
curved dominant-energy regions exist (Vaidya spacetimes before an incoming, or
after an outgoing, null-dust shell). Thin shells will not suffice in general,
since null junctions require matching intrinsic geometry. A solution would
give Theorem~\ref{thm:main} for spacetimes with physically reasonable matter;
a vacuum version would require bounded interior flat regions adjacent to
arbitrary curved vacuum regions, which we have not found in the gluing
literature (an abstract-level survey: Carlotto--Schoen, Corvino--Schoen,
Aretakis--Czimek--Rodnianski, Czimek--Rodnianski).


\subsection{What the classification shows}\label{sec:becoming-rival}

\paragraph{The robust part of Stein's theorem.} In general relativity, under
the premises of Remark~\ref{rem:lc}, nothing outside an event's causal past
can have become for it unless everything has. This holds in every globally
hyperbolic spacetime, however asymmetric, and does not rely on any foliation.
Those who want an extended, spacelike ``now'' as part of objective becoming
face Stein's dilemma in general relativity too, unless they change the
category or reject locality on the metric-only category.

\paragraph{The strongest rival: clock-equipped spacetimes.} Let $\GH_T$ have as
objects pairs $(M,T)$ with $M\in\GH$ and $T$ a smooth Cauchy temporal function, and
as morphisms $\psi:(M,T)\to(M',T')$ the morphisms of $\GH$ with $T'\circ\psi=T$. The
rule $x\,B_{(M,T)}\,b\iff T(x)\le T(b)$ is transitive, contains $\Ichr$, satisfies
(LC) on this category, and has the level sets of $T$ as shared presents. The
Milne inclusion of Proposition~\ref{prop:nonlocal} is not a morphism here: its
clock $\tau_W=(t^2-|\mathbf x|^2)^{1/2}$ has no smooth extension across the light
cone, where its gradient diverges. So this rival is local with respect to
its own morphisms. It rejects one premise, that becoming is a function of the
metric alone over all of $\GH$, and Theorem~\ref{thm:main} then says nothing about
it (in effect Monton's move). Its cost depends on the clock. A clock definable
from the metric, as cosmological time is in big-bang cosmologies, adds no
structure; the price is that the rule gives no verdict where that clock does
not exist, and locality no longer selects it, since $\Jm$ also satisfies the
axioms on $\GH_T$. A preferred foliation or reference field not definable
from the metric is added structure. Rejecting locality on $\GH$ itself is a
distinct option (Proposition~\ref{prop:nonlocal}).

\paragraph{A trilemma.} Suppose a shared present is understood as a set of
distinct events each of which has become for each other, the sense in which
Stein's dilemma is posed. Then Corollary~\ref{cor:nopresent} can be read as a
trilemma for becoming theories, which are transitive by definition, defined from the metric alone: locality, a
shared (extended) present, and non-trivial becoming cannot all be had. Defenders of a shared now reject locality on $\GH$, reject transitivity or change the category; local becoming theorists get
point-like presents, which recalls Stein's remark that the converse of
Putnam's premise leads to ``a peculiarly extreme (but pluralistic!) form of
solipsism'' \citep[p.~18]{Stein1968}; and giving up non-triviality is the
block. Other notions of a present, for instance a set of events none of which
has become for another, are not addressed by the corollary. The corollary
says nothing about which horn to take.

\paragraph{The fragile part.} Uniqueness does not survive. Minkowskian
symmetry forced a choice between $\Ichr$ and $\Jm$; general relativity leaves
open how lightlike-separated events stand, and a rule may let curvature along
the light ray decide. Clifton and Hogarth's debate over causal versus
chronological becoming \citep[pp.~364--365]{CH1995} thus acquires a third
family of options; the choice between the endpoints is equivalent to a
regularity requirement (Proposition~\ref{prop:reg}).

\paragraph{Locality does the work symmetry did.} In a single curved spacetime,
definability-style invariance constrains almost nothing
(Proposition~\ref{prop:single}); natural global rules such as cosmological
time violate only locality (Proposition~\ref{prop:nonlocal}). So the
curved-spacetime case for ``only the past has become'' rests on a locality
principle of the kind locally covariant field theory takes for granted and
akin to Viglione's Causal Connection Condition \citep{Viglione2026} (which is
stricter: (LC) allows dependence on spacelike events inside the convex region). Rejecting
locality is a coherent option, but it is the option of global, cosmological
becoming, with its own costs \citep{Dieks2006,Viglione2026}.

\paragraph{Limits.} The theorem concerns relations of becoming, not
ontology; it does not decide between presentism, the growing block and
eternalism. An eternalist can accept $\Jm$ as what has happened by an event.
Reading $B_M$ as existence from $b$ is a further bridge, which nothing here
supplies. The theorem assumes that becoming is defined for every globally hyperbolic
geometry (Remark~\ref{rem:offshell}); Section~\ref{sec:energy} shows that
dropping this premise is not routine. Spacetimes that are not globally hyperbolic are not treated beyond
Proposition~\ref{prop:godel}. The time orientation is an input, through (CB);
the theorem shows that, given an orientation, only the past lobe can have
become, but it does not explain the orientation. Nor does it address quantum
nonlocality. Clifton and Hogarth's maxim that whatever can causally influence
an event has become for it \citep[p.~364]{CH1995} needs care. Read as a claim
about signalling, it puts $\Jm$ into $B_M$. Bell-type dependence on spacelike
events is not signalling, so it is no counterexample to that reading. What it
challenges is the stronger thought that any dynamical dependence on an event
requires that event to have become. If the fundamental dynamics lets spacelike
events contribute to bringing an event about, the stronger reading puts them
into $B_M$, and Theorem~\ref{thm:main} makes becoming universal. A non-trivial
local becoming theory that holds the stronger reading therefore presupposes that
such contribution is restricted to the causal past. That is a bet about the
dynamics, not a theorem of relativity.


\subsection{Prior work on becoming}\label{sec:prior}

Stein's theorem and its Clifton--Hogarth form are the direct predecessors
\citep{Stein1968,Stein1991,CH1995}. \citet{Rakic1997} takes an event-relative,
transitive ``realized with respect to'' relation, argues that under her
constraints it is not definable from the causal geometry of Minkowski
spacetime, and adds it as a conservative primitive: in our terms, the
added-structure option. Point-based presents \citep{Muller2006}
and the causal-diamond present of Arthur and Savitt
\citep{Arthur2006,Arthur2019,Savitt2009,Savitt2015}, with Dorato's critique
\citep{Dorato2011} and Aames's application to Gödel universes \citep{Aames2022},
are conceptual proposals without a curved-spacetime classification.
\citet{Earman2008} builds worldline becoming in general relativity from causal
pasts and notes that once the symmetries of Minkowski spacetime are lost, the
Clifton--Hogarth definability results no longer apply \citep[p.~151]{Earman2008};
he does not pursue a definability result. In flat
spacetime, definability of sets from relativistic structure is governed by
automorphism groups \citep{MSS2025}.

The category and the naturality schema are those of \citet{BFV2003}; the
locality idea is theirs, and what is offered here is a classification of
relation-valued assignments. The closest structural precedent is the SPASs
property of \citet{FV2012}: theories conforming to a common notion of sameness,
one a subtheory of the other, that are isomorphic in one spacetime should be
isomorphic in all. Theorem~\ref{thm:main} has the same shape, one Minkowski verdict
fixing the alternative everywhere, but its objects are relations between events
and its mechanism is transitivity and (CB). Fewster and Verch show that the full
class of locally covariant theories fails the property, and
Theorem~\ref{thm:weyl} is a counterpart: $\Ichr$ and $\Ichr\cup N$ agree on
$\Mink$ and differ elsewhere. Stein himself posed a cross-model question of a
different kind, whether a cosmological model can be deformed to one with
Gödel-type relations \citep[p.~594]{Stein1970}. Clifton and Hogarth's remark
quoted after Corollary~\ref{cor:worldline} is an observation of difficulty, not a posed
problem. The flat-spacetime counterpart of Corollary~\ref{cor:nopresent}, that
relativity leaves no shared present under its symmetries, is established
\citep{Callender2000,Saunders2002}.

A citation sweep of
\citet{CH1995} and \citet{Stein1991} (Semantic Scholar, 24 October 2026; repeated
with OpenAlex, 25 October 2026), and
reading of the works named here, found no prior extension of
Stein's theorem to curved spacetime. This bounds, but does not certify,
priority.


A modal route to a relativistic present is related. M\"uller's branching-based proposal
\citep{Muller2019} defines co-presentness from branching structure, and its Theorem~1 asserts
non-trivial co-presentness on every history outside two featureless cases. Under the
necessary nonempty-chain reading, a comb with a dense set of branch points on a real trunk
satisfies that paper's Definition~6, is not featureless, and has identity co-presentness on
its trunk. In the comb family with branch set $D$ a proper subset of $\mathbb R$, all
histories have non-trivial co-presentness if and only if $D$ is not dense. That concerns an
abstract order, not a relativistic manifold, and the later formulation of the theory lies
outside the construction.

\partheading{Part II. Determination on a fixed background}

\section{What the causal structure guarantees}\label{sec:det}

Grant Part I. Does the unreality of $b$'s future from $b$, on views that read having
become as being real, make the future open from $b$? Openness in the relevant sense is a
matter of what the past fixes, and that depends on the laws and on global causal
structure. Geroch's domain of prediction supplies the tool \citep{Geroch1977}, and
Manchak's work on prediction its limits \citep{Manchak2008}.

\begin{definition}
An event $x$ is \emph{guaranteed from} $b$ if every past-inextendible causal
curve from $x$ meets $\Is(b)$. Write $\Det(b)$ for the set of such events.
Clearly $\Is(b)\subseteq\Det(b)$.
\end{definition}

$\Det(b)$ is a property of the spacetime and the event. To connect it with
what the laws fix, let $\mathcal L$ be the set of solutions of some field
equations on $M$. Say that $x$ is \emph{fixed from $b$ under $\mathcal L$},
and write $x\in F_{\mathcal L}(b)$, if any two solutions in $\mathcal L$ that
agree on $\Is(b)$ agree at $x$. The equations, background and globally
admissible solution class are held fixed in this comparison; they are inputs
to the definition, not data an observer is assumed to have measured in the
past. Here agreement means equality of the specified classical fields on the
whole open region, not equality of a finite collection of detector records.
Say that $\mathcal L$ has the
\emph{domain-of-dependence property} if, for every open $U\subseteq M$, any two
solutions that agree on $U$ agree at every event all of whose
past-inextendible causal curves meet $U$. We assume this uniqueness property
explicitly; it is not inferred merely from an operational no-signalling
condition.

\begin{remark}[What is guaranteed and what is fixed]\label{rem:laws}
If $\mathcal L$ has the domain-of-dependence property and consists of
continuous fields, then
\[\Det(b)\cup J^-(b)\subseteq F_{\mathcal L}(b),\]
since $J^-(b)$ is the closure of $\Is(b)$ in a globally hyperbolic spacetime
\citep{MS2008}. So $\Det(b)$ is fixed under every such law: it is what the
causal structure guarantees, whatever the fields. The inclusion is in general
strict, and nothing about $F_{\mathcal L}(b)$ beyond it follows from $\Det(b)$
being small. For example, on Minkowski spacetime with the unit timelike
field $u=\partial_t$, the local equations $u\cdot\nabla Q=P$,
$u\cdot\nabla P=0$ have the solutions $Q=A(\vec X)+tB(\vec X)$, $P=B(\vec X)$
for arbitrary smooth $A,B$. The values of $(Q,P)$ at one point of an integral curve fix them along
it, and the past integral curve from $x$ is a past-inextendible causal curve,
so the law has the domain-of-dependence property; and every integral curve
of $u$ enters $\Is(b)$, so $F_{\mathcal L}(b)=M$ for every $b$. A law of this
kind uses background structure beyond the metric, but nothing in the
domain-of-dependence property excludes it. Neither notion is Belnap's
settledness, on which a proposition is settled at a point event if it holds
in every possible history through it \citep{Belnap1992}; under
indeterministic laws they come apart from it, and we say nothing about such
laws. We use $\Is(b)$, which every local becoming theory counts as real from
$b$, rather than $J^-(b)$; the characterization below uses that $\Is(b)$ is
open.
\end{remark}

\paragraph{Relation to Geroch.} Geroch's domain of prediction $P(q)$ is the set
of $x$ such that every past-directed timelike curve from $x$ without past
endpoint enters $I^-(q)$, but $I^-(x)\not\subseteq I^-(q)$ \citep[sec.~3]{Geroch1977}.
$\Det(b)$ is the causal-curve analogue of the union of $P(b)$ with the
events whose pasts lie in that of $b$; we expect the two to differ at most
on null boundaries, but have not proved it. Geroch gives an equivalent form by achronal surfaces in $I^-(q)$
\citep[sec.~3]{Geroch1977}; Theorem~\ref{thm:char} is a slicing form for
globally hyperbolic spacetimes.

\begin{theorem}\label{thm:char}
Let $t:M\to(\alpha,\beta)$ be a Cauchy time function. Then $x\in\Det(b)$ if
and only if there is $c\in(\alpha,t(x))$ with $J^-(x)\cap S_c\subseteq\Is(b)$.
\end{theorem}

\begin{proof}
($\Leftarrow$) A past-inextendible causal curve from $x$ crosses the Cauchy
surface $S_c$, at a point of $J^-(x)\cap S_c\subseteq\Is(b)$.

($\Rightarrow$) Suppose that for every $c\in(\alpha,t(x))$ the set
$K_c=(J^-(x)\cap S_c)\setminus\Is(b)$ is nonempty; it is compact, since
$\Is(b)$ is open. Fix $t(x)>c_1>c_2>\dots$ with $c_n\to\alpha$, and let
\[P_n=\{(z_1,\dots,z_n)\in K_{c_1}\times\dots\times K_{c_n}:
x\ge z_1\ge z_2\ge\dots\ge z_n\}.\]
$P_n$ is compact, since $\le$ is closed. It is nonempty: pick $y\in K_{c_n}$
and a causal curve from $x$ to $y$; it crosses each $S_{c_i}$, $i<n$, at some
$z_i\le x$, and $z_i\notin\Is(b)$, for otherwise $y\le z_i\ll b$ would give
$y\ll b$. The maps $P_{n+1}\to P_n$ forgetting the last coordinate are
continuous, so the inverse limit of the nonempty compact spaces $P_n$ is
nonempty, giving $x\ge z_1\ge z_2\ge\cdots$ with $z_n\in K_{c_n}$. Join
consecutive terms by causal curves and concatenate. Each segment avoids
$\Is(b)$ by push-up, as before. Along the concatenation $t\to\alpha$, so it has
no past endpoint in $M$: it is a past-inextendible causal curve from $x$
missing $\Is(b)$.
\end{proof}

\begin{corollary}\label{cor:cauchy}
If $\Is(b)$ contains a Cauchy surface $S$, then $\Det(b)=M$.
\end{corollary}

\begin{proof}
If $x\in\Is(S)\cup S$, then $x\le s\ll b$ for some $s\in S$, so $x\in\Is(b)$.
If $x\in\Isp(S)$, extend a past-inextendible causal curve from $x$ to an
inextendible one; it meets $S$, and not to the future of $x$, since then
$s'\ll x\le s$ for some $s,s'\in S$, contradicting achronality. So the curve
from $x$ meets $S\subseteq\Is(b)$.
\end{proof}

Since $\Det(b)$ is defined by causal curves alone, it depends only on the
conformal class of the metric; conformal diagrams decide it.


\section{The standard models}\label{sec:examples}

\begin{lemma}\label{lem:balls}
Let $(\Sigma,d)$ be $\mathbb R^3$, hyperbolic space $H^3$, or the unit sphere
$S^3$ with $r_2\le\pi$. Then the closed ball $\bar B(X,r_1)$ lies in the open
ball $B(Y,r_2)$ if and only if $d(X,Y)+r_1<r_2$.
\end{lemma}

\begin{proof}
($\Leftarrow$) is the triangle inequality. ($\Rightarrow$) Let $r_1\ge0$ and suppose
$d(X,Y)+r_1\ge r_2$. On $\mathbb R^3$ and $H^3$, where geodesics extend
indefinitely and minimize, or on $S^3$ when $d(X,Y)+r_1\le\pi$, follow a minimizing geodesic from $Y$ through $X$ (any geodesic from
$X$ if $X=Y$) a further distance $r_1$; the endpoint lies in $\bar B(X,r_1)$ at
distance $d(X,Y)+r_1\ge r_2$ from $Y$. If $d(X,Y)+r_1>\pi$ on $S^3$, the
antipode of $Y$ is at distance $\pi-d(X,Y)<r_1$ from $X$ and $\pi\ge r_2$ from
$Y$.
\end{proof}

\begin{proposition}\label{prop:conformal}
Let $M$ be conformal to $(\eta_-,\eta_+)\times\Sigma$ with metric
$-d\eta^2+d\Sigma^2$, where $\Sigma$ is $\mathbb R^3$, $H^3$, or the unit $S^3$
with $\eta_+-\eta_-\le\pi$. Then $\Det(b)=\Is(b)$ for every $b$.
\end{proposition}

\begin{proof}
Use $t=\eta$. For $x=(\eta_x,X)$, $b=(\eta_b,Y)$: $x\ll b$ iff
$d(X,Y)<\eta_b-\eta_x$ (on $S^3$, $\eta_b-\eta_x<\pi$). For $\eta_-<c<\eta_x$,
$J^-(x)\cap S_c=\bar B(X,\eta_x-c)$ and $\Is(b)\cap S_c=B(Y,\eta_b-c)$, with
$\eta_b-c<\pi$ on $S^3$. By Lemma~\ref{lem:balls} the inclusion holds iff
$d(X,Y)<\eta_b-\eta_x$, i.e.\ iff $x\ll b$. Apply Theorem~\ref{thm:char}.
\end{proof}

This covers Minkowski spacetime, spatially flat and open FLRW universes with
a big bang (conformal to $\eta>0$), the closed radiation universe
($a\propto\sin\eta$, $\eta\in(0,\pi)$) and global de Sitter spacetime ($\Sigma=S^3$,
$\eta\in(-\pi/2,\pi/2)$) \citep{HE1973}. In none of them is anything outside an
event's chronological past guaranteed from it. In Minkowski spacetime and the
spatially flat universes this bound is attained by a natural field.

\begin{proposition}\label{prop:sharp}
Let $M=(\eta_-,\eta_+)\times\mathbb R^3$ with $g=a(\eta)^2(-d\eta^2+d\vec X^2)$,
$a>0$ smooth, and let $\mathcal L$ be the smooth solutions of the conformally
coupled scalar equation $\Box_g\phi-\tfrac16R\phi=0$. Then
$F_{\mathcal L}(b)=J^-(b)$ for every $b$.
\end{proposition}

\begin{proof}
$J^-(b)\subseteq F_{\mathcal L}(b)$ by continuity, as in
Remark~\ref{rem:laws}. Conversely, $\phi$ is a solution iff $\psi=a\phi$
solves the flat wave equation $\partial_\eta^2\psi=\Delta\psi$ on $M$
\citep[app.~D]{Wald1984}. Let $b=(\eta_b,\vec Y)$ and $x=(\eta_x,\vec X)\notin
J^-(b)$, so $\eta_x-\eta_b+|\vec X-\vec Y|>0$, since the causal curves of $M$
are those of the conformal Minkowski metric and the slab is causally convex in
Minkowski spacetime. Choose a unit vector $\hat k$
with $\hat k\cdot(\vec X-\vec Y)=-|\vec X-\vec Y|$, a smooth $f$ vanishing on
$(-\infty,0]$ and positive on $(0,\infty)$, and put $\psi=f(u)$ with
$u=\eta-\eta_b-\hat k\cdot(\vec Z-\vec Y)$ at $(\eta,\vec Z)$. This solves the
wave equation. On $\Is(b)$, $\eta-\eta_b<-|\vec Z-\vec Y|$, so $u<0$ and
$\psi=0$; at $x$, $u>0$ and $\psi\neq0$. So $\psi/a$ and $0$ are solutions that
agree on $\Is(b)$ but not at $x$.
\end{proof}


\begin{proposition}\label{prop:dust}
Let $M$ be conformal to $(\eta_-,\eta_+)\times S^3$ with $\eta_+-\eta_->\pi$.
Then $\Det(b)=\Is(b)$ if $\eta_b\le\eta_-+\pi$, and $\Det(b)=M$ if
$\eta_b>\eta_-+\pi$. In the Einstein static universe ($\eta_\pm=\pm\infty$),
$\Det(b)=M$ for every $b$.
\end{proposition}

\begin{proof}
If $\eta_b>\eta_-+\pi$, take $c\in(\eta_-,\eta_b-\pi)$: $\Is(b)\cap S_c$ is an open
ball of radius greater than $\pi$, hence all of $S_c$, and
Corollary~\ref{cor:cauchy} applies. If $\eta_b\le\eta_-+\pi$, every $c>\eta_-$
has $\eta_b-c<\pi$, and the proof of Proposition~\ref{prop:conformal} applies.
\end{proof}

For the closed dust FLRW model, $a(\eta)=\tfrac12a_m(1-\cos\eta)$ with
$\eta\in(0,2\pi)$, the threshold is $\eta=\pi$, the moment of maximum
expansion. That the particle horizon of the closed dust model disappears at
maximum expansion, so that an observer can then in principle see the whole
universe, is known \citep{GronJohannesen2011,BolotinTanatarov2013}; the
Einstein static case is Geroch's \citep[Fig.~5]{Geroch1977}. What
Proposition~\ref{prop:dust} adds is the consequence for determination: with
Corollary~\ref{cor:cauchy}, once the particle horizon has vanished
($\eta_b>\pi$), everything is guaranteed from $b$, so that every law with the
domain-of-dependence property fixes everything from what is real from $b$.
This fits Manchak's theorem that if some event of $b$'s timelike future lies
in Geroch's domain of prediction of $b$, the spacetime admits a compact
spacelike slice \citep[p.~319]{Manchak2008}. That theorem concerns the
geometric notion. The flow law of Remark~\ref{rem:laws} shows that it has no
analogue for what a particular law fixes: in Minkowski spacetime, which has no
compact slice, that law fixes all of $b$'s timelike future. Manchak also
records the view that ``Maxwell's equations constrain some spacetimes in such
a way as to allow for prediction'' \citep[p.~321 and n.~15]{Manchak2008}, so
the law-relative notion has antecedents. Earman's discussion explicitly
includes the Cauchy-surface exception and an electromagnetic-constraint
example citing Ellis and Sciama \citep[pp.~125--130]{Earman1995}. We have checked
these selected original Earman pages; Ellis and Sciama (1972, pp.~41--42)
remain unread here and enter this comparison through Earman's discussion.
The model-specific characterization below may overlap with those sources;
we make no priority claim for it.

Before maximum expansion a particular field can fix more than is guaranteed.
Light leaving a point of the closed model refocuses at its antipode after
conformal time $\pi$, and the rescaled conformal scalar $\psi=a\phi$ returns
there with its sign reversed.

\begin{proposition}\label{prop:echo}
Let $M$ be the closed dust model, $g=a(\eta)^2(-d\eta^2+d\Omega_3^2)$ with
$a=\tfrac12a_m(1-\cos\eta)$, $\eta\in(0,2\pi)$, and $\mathcal L$ as in
Proposition~\ref{prop:sharp}. Let $A(\eta,X)=(\eta+\pi,-X)$, where $-X$ is the
antipode of $X$. Then $J^-(b)\cup\bigl(A(J^-(b))\cap M\bigr)\subseteq
F_{\mathcal L}(b)$. The set $A(J^-(b))\cap M$ is disjoint from $\Isp(b)$, and if
$\eta_b\le\pi$ it is nonempty and disjoint from $\Det(b)$.
\end{proposition}

\begin{proof}
Here $\psi=a\phi$ solves the conformally invariant equation of the Einstein
static universe, $\partial_\eta^2\psi=\Delta\psi-\psi$ on the unit $S^3$
\citep[app.~D]{Wald1984}, and by global well-posedness extends uniquely to
$\mathbb R\times S^3$.
Expanding in spherical harmonics of degree $n$, which have frequency $n+1$
and parity $(-1)^n$ under the antipodal map, each mode, and so every smooth
solution, satisfies $\psi(\eta+\pi,-X)=-\psi(\eta,X)$. Two solutions that agree on $\Is(b)$ agree, by continuity, on $J^-(b)$, and
hence on $A(J^-(b))\cap M$. If $(\eta,X)\in J^-(b)$ and $\eta+\pi<2\pi$, then
$d(X,Y)\le\eta_b-\eta$, so $d(-X,Y)=\pi-d(X,Y)\ge(\eta+\pi)-\eta_b$;
thus $A(\eta,X)\notin\Isp(b)$. If $\eta_b\le\pi$, then
$\Det(b)=\Is(b)$ by Proposition~\ref{prop:dust}, which lies in
$\{\eta<\eta_b\}$, while $A(J^-(b))$ lies in $\{\eta>\pi\}$ and meets $M$, for
instance at $A(\eta_b/2,Y)$ when $b=(\eta_b,Y)$.
\end{proof}

For $\eta_b\le\pi$, the echo events are unreal from $b$ and fixed from $b$,
but they lie beside $b$ or on its future light cone, not in its timelike
future. The field fixes
more than the echo, including part of $b$'s timelike future. Let
$\Diamond_b=\{(\eta,X)\in\mathbb R\times S^3:d(X,Y)+|\eta|<\eta_b\}$, the double
cone in the Einstein static universe between $b$ and its mirror image
$(-\eta_b,Y)$, and $U_b=\Diamond_b\cup A\Diamond_b\cup A^2\Diamond_b$.

\begin{theorem}\label{thm:diamonds}
In the setting of Proposition~\ref{prop:echo}, let $0<\eta_b<\pi$. Then
$F_{\mathcal L}(b)=M\cap\overline{U_b}$. In
particular $F_{\mathcal L}(b)$ contains the nonempty set
$\{(\eta,X)\in M:d(X,Y)\le\eta+\eta_b-2\pi\}$, which lies in $\Isp(b)$.
\end{theorem}

\begin{proof}
For data $(f,g)$ at $\eta=0$, the Einstein static solution is
$\psi(\eta,X)=\partial_\eta[\sin\eta\,M_f(X,\eta)]+\sin\eta\,M_g(X,\eta)$, where
$M_h(X,r)$ is the mean of $h$ over the geodesic sphere of radius $r$ about
$X$; one checks this on each mode, using that the zonal function of degree $n$
on $S^3$ is $\sin((n+1)r)/((n+1)\sin r)$. So $\psi(x)$ depends on the data only
near the sphere $\Sigma_x=S(X,\eta_x)$, which is the point $-X$ when $\eta_x=\pi$.

\emph{Sufficiency.} Solutions agreeing on $\Is(b)$ agree on the ball
$B(Y,\eta_b-\epsilon)$ in the slice $\eta=\epsilon$, hence on its domain of
dependence in the Einstein static universe; the union over $\epsilon>0$ is
$\Diamond_b$. The identities $\psi\circ A=-\psi$ and $\psi\circ A^2=\psi$ carry
the agreement to $A\Diamond_b$ and $A^2\Diamond_b$.
Moreover, $F_{\mathcal L}(b)$ is relatively closed in $M$: it is the
intersection of the zero sets of all smooth difference solutions vanishing
on $\Is(b)$. Since $M$ is open in the Einstein cylinder, taking the relative
closure gives $M\cap\overline{U_b}\subseteq F_{\mathcal L}(b)$.

\emph{Necessity.} Let $x\in M\setminus\overline{U_b}$. Replacing $x$ by
$A^{-1}x$ if $\eta_x>\pi$, assume $\eta_x\le\pi$. Let $C=\bar B(Y,\eta_b)$.
Data supported off $C$ give a solution vanishing on $\Is(b)$: if $z\in\Is(b)$
lay in its support, then $w\le z\ll b$ for some $w$ in the support of the data,
so $d(w,Y)<\eta_b$. If $\eta_x=\pi$, then $d(-X,Y)>\eta_b$, and data
$(f,0)$ with $f$ a bump at $-X$ off $C$ give $\psi(x)=-f(-X)\neq0$. If
$\eta_x<\pi$, then $\Sigma_x\not\subseteq C$: otherwise the connected set
$S^3\setminus C$ lies in one of the two open balls bounded by $\Sigma_x$, so
the other closed ball lies in $C$, which by the argument of
Lemma~\ref{lem:balls} means
$d(X,Y)+\eta_x\le\eta_b$ or $d(-X,Y)+\pi-\eta_x\le\eta_b$, that is,
$x\in\overline{\Diamond_b}\cup\overline{A\Diamond_b}$. Take $p\in\Sigma_x\setminus C$
and data $(0,g)$ with $g\ge0$ supported near $p$ off $C$ and $g(p)>0$; then
$\psi(x)=\sin\eta_x\,M_g(X,\eta_x)>0$.

Finally, if $d(X,Y)\le\eta+\eta_b-2\pi$, then $d(X,Y)<\eta-\eta_b$ because
$\eta_b<\pi$, so $(\eta,X)\in\Isp(b)$; the set contains $(\eta,Y)$ for
$\eta\ge2\pi-\eta_b$.
\end{proof}

On $b$'s comoving geodesic $X=Y$, the field is fixed from $b$ exactly for
$\eta\le\eta_b$ and for $\eta\ge2\pi-\eta_b$, and open in between: the end is
fixed before the middle. The reason is that the past light cone of an event
just before the crunch wraps around the universe and refocuses near $Y$ at the
big bang, inside what is real from $b$. The same sphere argument shows that,
up to boundary, $M\cap U_b$ is the set
of events all of whose past-inextendible null geodesics meet $\Is(b)$, so for
this field, which propagates on light rays alone, the null geodesics decide
what is fixed. The same argument applies to global de Sitter spacetime,
conformal to the slab $(-\pi/2,\pi/2)$, where the field is fixed from $b$ on the
antipodal comoving geodesic for all $\eta>-\eta_b$, although no signal from that
geodesic ever reaches the comoving geodesic of $b$. So the geometric threshold
at maximum expansion marks when everything is guaranteed, not when this field
first fixes anything beyond the past.

For a law whose fixed set itself changes at the threshold, take the comoving
flow $u=a^{-1}\partial_\eta$ of the closed dust model with the equations of
Remark~\ref{rem:laws}; here $u$ is the dust's own four-velocity, so the law
adds no structure the model lacks. The comoving line through $X$ meets $\Is(b)$ iff
$d(X,Y)<\eta_b$, so, by continuity in $X$ and by choosing data supported off
the closed ball, $F_{\mathcal L}(b)=(0,2\pi)\times\bar B(Y,\eta_b)$. This is a
proper subset of $M$, containing events in $b$'s timelike future, if
$\eta_b<\pi$, and all of $M$ if $\eta_b\ge\pi$. Under this law the answer to
``is the future fixed from here?'' differs between events before and after
maximum expansion.

Theorem~\ref{thm:massive} shows that on closed dust a local law can fix exactly
$\Det(b)\cup J^-(b)$ for every $\tau\neq\pi$. Whether some local law does so at
$\tau=\pi$, or in other closed models, we leave open.


\section{Changing the law}\label{sec:mass}

Write $\tau=\eta_b$ for the conformal time of a standpoint.


The massless example is special. Its extra fixing can change even when the
metric and its causal curves are held fixed. Let $\mathcal L_m$ consist of
all smooth real solutions on the closed dust strip of
\[
 (\Box_g-R_g/6-m^2)\phi=0,
\]
with ordinary constant physical mass $m\ge0$, and write $F_m(b)$ for its
fixing set. Set $\Lambda^2=-\Delta_{S^3}+1$ and $T=2\pi$. With
$\psi=a\phi$ and $\lambda=m^2$, the rescaled equation is
\begin{equation}\label{eq:mass}
 \psi_{\lambda,\eta\eta}+(\Lambda^2+\lambda q(\eta))\psi_\lambda=0,
 \qquad q(\eta)=a(\eta)^2.
\end{equation}
In particular a constant addition to the rescaled potential would be a
different perturbation, corresponding to physical potential $\delta/a^2$.

\begin{proposition}[Loss of exact fixing under ordinary mass]\label{prop:mass}
Fix the dust amplitude $a_m>0$. For each $b=(\tau,Y)$ with $0<\tau<\pi$,
there are a fixed event
$x_*=(T-s,Y)$, $0<s<\tau$, and $m_*>0$ such that
\[
 x_*\in F_0(b)\cap\Isp(b),\qquad
 x_*\notin F_m(b)\quad\text{for every }0<m<m_*.
\]
The background metric is the same for all these solution classes.
\end{proposition}

\begin{proof}
Choose a nonzero smooth $v\ge0$ supported strictly outside
$\bar B(Y,\tau)$, and solve \eqref{eq:mass} with data $(0,v)$ at the
auxiliary cylinder slice $\eta=0$. The coefficient $q=a^2$ is smooth through
$0$ and $T$, so ordinary smooth Cauchy evolution and parameter dependence
apply on the extended cylinder. This is a device for constructing solutions
on the open physical strip, not a physical continuation through its
singularities. For $z=(\eta,Z)\in\Is(b)$, the initial ball
$\bar B(Z,\eta)$ lies in $B(Y,\tau)$ by the triangle inequality. Finite
propagation therefore gives $\psi_\lambda=0$ throughout $\Is(b)$ for every
$\lambda$; the zero-order term does not change this conclusion.

Put $S_0(u)=\sin(u\Lambda)/\Lambda$. The first variation
$w=\partial_\lambda\psi_\lambda|_{\lambda=0}$ has zero initial data and source
$-q(u)S_0(u)v$. Since $S_0(T-u)=-S_0(u)$, Duhamel's formula gives
\[
 w(T)=\int_0^T q(u)S_0(u)^2v\,du.
\]
The frequencies are $k=n+1$. Expanding
$q=(a_m^2/4)(3/2-2\cos u+\tfrac12\cos2u)$ gives
\[
 \int_0^T q(u)\frac{\sin^2(ku)}{k^2}\,du
 =\begin{cases}5\pi a_m^2/16,&k=1,\\
                 3\pi a_m^2/(8k^2),&k\ge2.
   \end{cases}
\]
Thus, with $\Pi_0$ the projection onto spatial constants,
\[
 w(T)=\left(\frac{3\pi a_m^2}{8}\Lambda^{-2}
                -\frac{\pi a_m^2}{16}\Pi_0\right)v.
\]
The kernel of $\Lambda^{-2}$ on the unit $S^3$ is
$G(r)=(\pi-r)/(4\pi^2\sin r)$: it solves the homogeneous radial equation
away from its pole, is smooth at the antipode, and has the normalized
singularity $1/(4\pi r)$. Since $\Pi_0$ has kernel $1/(2\pi^2)$, the kernel
of the operator displayed above is
\[
 K(r)=\frac{a_m^2}{32\pi}
       \left(\frac{3(\pi-r)}{\sin r}-1\right)>0\qquad(0<r\le\pi),
\]
using $(\pi-r)/\sin r\ge1$, with its continuous value at $r=\pi$.
Consequently $w(T,Y)>0$. Continuity now permits a choice of a fixed
$s\in(0,\tau)$ with $w(T-s,Y)>0$. Theorem~\ref{thm:diamonds} puts
$x_*=(T-s,Y)$ in $F_0(b)\cap\Isp(b)$, so $\psi_0(x_*)=0$.
Smooth dependence on $\lambda$ yields
\[
 \psi_{m^2}(x_*)=m^2w(x_*)+o(m^2)>0
\]
for every sufficiently small positive $m$. Because $a(x_*)>0$, the physical
field and zero agree on $\Is(b)$ but differ at $x_*$. The event and data
were chosen before the mass interval.
\end{proof}

Theorem~\ref{thm:massive} below gives the full formula; this proposition
supplies an explicit first-order witness. The result concerns a
linear classical field on a fixed background, without interactions,
quantization or gravitational backreaction.


\begin{theorem}[Massive fields, complete classification]\label{thm:massive}
On the closed dust model, for each $m>0$ let $\mathcal L_m$ be all smooth real solutions of
$(\Box_g-R_g/6-m^2)\phi=0$, and let $b=(\tau,Y)$. Then
\[
F_m(b)=\Jm(b)\quad\text{if }0<\tau<\pi,\qquad F_m(b)=M\quad\text{if }\pi\le\tau<2\pi.
\]
The second line also holds for $m=0$.
\end{theorem}

\begin{proof}[Outline]
Put $\psi=a\phi$, so that $\psi_{\eta\eta}+(\Lambda^2+V(\eta))\psi=0$ with $V=m^2a^2$ entire and positive
on $(0,2\pi)$. A successive-approximation construction gives a line Riemann function jointly
analytic in its variables. Periodizing it gives the displacement and velocity kernels on
$S^3$. The analytic pieces of the kernels are not identically zero, so the joint support of
the two Cauchy kernels has no open holes. For $\tau<\pi$ and $x\notin\Jm(b)$, the kernel
support therefore leaves the closed ball $\bar B(Y,\tau)$, and data supported off that ball
give a solution vanishing on $\Is(b)$ and not at $x$; the reverse inclusion is continuity.
For $\tau>\pi$ the trace of $\Is(b)$ on an early slice is the whole sphere. For $\tau=\pi$ it is
the sphere minus a point, and smoothness forces both data to vanish.
The full argument is in the project's public note (Appendix~\ref{app:status}).
\end{proof}

So before maximum expansion any positive mass removes every island of extra fixing, while
the threshold itself does not move, and at the exact threshold $\tau=\pi$ the law fixes
everything although $\Det(b)=\Is(b)$. Dense traces do this work there, not causal curves.

\subsection{A norm bound for hidden alternatives}

Exact fixing quantifies over every admissible difference, however small.
To ask about its size, a restriction on alternatives must be supplied. Use
the conformal Cauchy-data norm
\[
 N(f,v)^2=\|\Lambda f\|_2^2+\|v\|_2^2
\]
at the auxiliary slice $\eta=0$, with unit-$S^3$ volume $2\pi^2$.
This is not an observer's measured physical-energy budget.

\begin{proposition}\label{prop:norm}
Let $\psi_\lambda$ and $\psi_0$ have the same smooth data $(f,v)$ and solve
the cylinder equation with spatially homogeneous smooth real potential
$\lambda q(\eta)$ and with zero potential, respectively. For $\eta\ge0$, put
$Q(\eta)=\int_0^\eta |q(u)|\,du$. Then
\begin{equation}\label{eq:norm}
 \sup_X|\psi_\lambda(\eta,X)-\psi_0(\eta,X)|
 \le \frac{|\lambda|Q(\eta)e^{|\lambda|Q(\eta)}}{\sqrt6}\,N(f,v).
\end{equation}
If $f=0$, the denominator improves to $\sqrt{12}$.
\end{proposition}

\begin{proof}
Let $C_{\lambda,k}$ and $S_{\lambda,k}$ be the fundamental solutions with
data $(1,0)$ and $(0,1)$ at frequency $k\ge1$. Their Volterra equations and
Gronwall's inequality imply
$|C_{\lambda,k}|\le e^{|\lambda|Q/k}$ and
$|S_{\lambda,k}|\le e^{|\lambda|Q/k}/k$. Substitution back into those
equations bounds the differences from the free multipliers by
\[
 |C_{\lambda,k}-\cos k\eta|\le B/k,\qquad
 |S_{\lambda,k}-\sin(k\eta)/k|\le B/k^2,
 \quad B=|\lambda|Qe^{|\lambda|Q}.
\]
Acting on $\Lambda f$ rather than $f$, the first multiplier acquires another
$1/k$. The addition identity for an orthonormal harmonic basis is
$\sum_j|Y_{kj}(X)|^2=k^2/(2\pi^2)$, since the eigenspace has dimension
$k^2$ and its diagonal projector is rotation-invariant. Joint
Cauchy--Schwarz on the coefficients of $(\Lambda f,v)$ bounds the squared
evaluation norm by
\[
 \sum_{k=1}^\infty\frac{k^2}{2\pi^2}
       \left(\frac{B^2}{k^4}+\frac{B^2}{k^4}\right)
 =\frac{B^2}{6}.
\]
Removing the position-data term halves this bound. The argument concerns
smooth solutions and the bounded difference functional; it does not assign
point values to arbitrary individual finite-energy fields.
\end{proof}

Apply this to any rescaled difference $h_m=a(\phi_1-\phi_2)$ of two massive
solutions agreeing exactly on $\Is(b)$, with data $(f,v)$. Smooth extension to zero
implies that both data vanish on $B(Y,\tau)$. The massless solution $h_0$
with the same data vanishes on $\Is(b)$ by finite propagation. Thus
$h_0(x)=0$ for every $x\in F_0(b)$, and \eqref{eq:norm} gives
\begin{equation}\label{eq:hidden}
 |h_m(\eta,X)|\le
 \frac{m^2Q(\eta)e^{m^2Q(\eta)}}{\sqrt6}N(f,v),
 \qquad (\eta,X)\in F_0(b),
\end{equation}
where, for dust,
\[
 Q(\eta)=\frac{a_m^2}{4}
 \left(\frac{3\eta}{2}-2\sin\eta+\frac{\sin2\eta}{4}\right).
\]
The bound controls all such smooth alternatives with $N\le1$, not just the
witness in Proposition~\ref{prop:mass}. At fixed targets their rescaled
ambiguity is $O(m^2)$. For physical fields divide by $a(\eta)>0$; uniform
physical-field control requires staying away from the singular endpoints.
Without a norm bound any nonzero hidden solution can instead be scaled to
an arbitrarily large target difference. Neither convention is an observation
model supplied by the geometry.

\subsection{Errors in records are a different question}

Equation~\eqref{eq:hidden} retains exact agreement on the whole past while
changing the law. It is not a stability theorem for noisy records under a
fixed law. Even at a massless-fixed target, an aggregate field-error norm
need not control a point value. To see the regularity issue explicitly, fix
any interior target $(\eta_*,Y)$ and let $Z_k$ be the real unit-$L^2$ zonal
harmonic centered at $Y$, with $Z_k(Y)=k/\sqrt{2\pi^2}$. The smooth massless
solutions
\[
 \psi_K(\eta,X)=\frac1{\sqrt K}
  \sum_{k=K}^{2K-1}\frac{\cos(k(\eta-\eta_*))}{k}Z_k(X)
\]
have conserved conformal norm $N=1$ and target value
$\sqrt{K/(2\pi^2)}$, but
\[
 \|\psi_K\|_{L^2((0,2\pi)\times S^3)}^2
 =\frac\pi K\sum_{k=K}^{2K-1}\frac1{k^2}\le\frac\pi{K^2}.
\]
Their error on $\Is(b)$ is even smaller. They do not agree \emph{exactly}
there. This is unbounded point evaluation under a specified norm, already
present with records of the entire strip, not a causal-horizon or detector
failure. Every fixed spatial average with profile $\chi\in L^2(S^3)$ instead
satisfies $|\langle\chi,\psi_K(\eta)\rangle|\le\|\chi\|_2/K$.
No general past-record stability theorem for such averages is asserted.

A common finite spectral cutoff changes the exact question too: every
solution in that restricted massless class is a finite real-analytic sum,
so agreement on any nonempty open region entails agreement everywhere.
Its fixing set is all of $M$. Thus finite mode sums cannot establish exact
past agreement with a nonzero future difference in the unrestricted smooth
theory. These are classical regularity observations, not further priority
claims or results about the physical measurement process.


\section{The fixing map}\label{sec:map}

The results combine into a single picture. For a law class $\mathcal L$ of continuous fields
with the domain-of-dependence property, and a non-universal becoming theory,
\begin{equation}\label{eq:chain}
\Is(b)\ \subseteq\ B(b)\ \subseteq\ \Jm(b)\ \subseteq\ F_{\mathcal L}(\Is(b)),
\qquad \Is(b)\ \subseteq\ \Det(b)\ \subseteq\ F_{\mathcal L}(\Is(b)).
\end{equation}

\begin{theorem}[All patterns occur]\label{thm:map}
Each row below holds, and each of $\Jm(b)\subseteq F$, $\Is(b)\subseteq\Det(b)$ and $\Det(b)\subseteq F$ can be strict or an equality. Here $B(b)=\{x:x\,B_M\,b\}$ and $F=F_{\mathcal L}(\Is(b))$.
\begin{center}
\begin{tabularx}{\textwidth}{@{}Yll@{}}
\toprule
Spacetime, law, standpoint & $\Det(b)$ & $F(b)$\\
\midrule
Minkowski, conformal scalar & $\Is(b)$ & $\Jm(b)$\\
Minkowski, preferred-flow law & $\Is(b)$ & $M$\\
Closed dust, $\tau<\pi$, conformal scalar & $\Is(b)$ & $M\cap\overline{U_b}$, three diamonds\\
Closed dust, $\tau<\pi$, comoving flow law & $\Is(b)$ & $(0,2\pi)\times\bar B(Y,\tau)$\\
Closed dust, $\tau<\pi$, mass $m>0$ & $\Is(b)$ & $\Jm(b)$\\
Closed dust, $\tau=\pi$, mass $m\ge0$ & $\Is(b)$ & $M$\\
Closed dust, $\tau>\pi$, every law above & $M$ & $M$\\
\bottomrule
\end{tabularx}
\end{center}
\end{theorem}

\begin{proof}
\eqref{eq:chain} is Remark~\ref{rem:laws} with Theorem~\ref{thm:main}. Row one is Propositions~\ref{prop:conformal} and~\ref{prop:sharp}; row two Proposition~\ref{prop:conformal} and Remark~\ref{rem:laws}; rows three
and four Proposition~\ref{prop:dust}, Theorem~\ref{thm:diamonds} and the flow calculation after it;
rows five and six Proposition~\ref{prop:dust} and Theorem~\ref{thm:massive}; row seven
Proposition~\ref{prop:dust} with Corollary~\ref{cor:cauchy}.
\end{proof}

Read down the table. Reading having become as being real, being unreal from $b$ does not imply being open: rows two to four, six and
seven. Being guaranteed is a lower bound that says nothing about what a law adds: rows two to four and six. The same metric gives different answers under different laws: rows three to five.
Determinism does not make the future fixed from a standpoint: in row one the law has unique
Cauchy evolution and nothing beyond the past is fixed. And an exact verdict can change while the
fields barely do: by \eqref{eq:hidden}, the differences hidden at a massless-fixed target are
$O(m^2)$ in a stated norm. Exact nonfixing and practical unpredictability are separate questions.

\partheading{Part III. Determination with gravity}

\section{Same past, different futures}\label{sec:coupled}

A fixed background leaves room for a reply. Perhaps gravity ties things together: perhaps
coupling, conserved mass or regularity at the center would make the past of an event determine
what later happens there. Subject to the inputs below, it does not.

Work with the spherically symmetric Einstein--real massless scalar system. Use double null coordinates $(u,v)$ with $t=(u+v)/2$ and a launch event $b$ at $t=0$ on the center, with future unit vector $n$ along it; $r$ is the area radius and $m$ the Hawking mass. Choose incoming characteristic data at past null infinity from a compact family: fix $0<A<T_1$, smooth $F_k$ with disjoint supports about $T_k$ inside $(A,2T_N+1)$ and $F_k(T_j)=\delta_{kj}$, and a smooth nonzero $G$ supported after $2T_N+1$. Take $x$ in a small closed ball on which $\beta$ below is smooth and positive. Set
\[
\Phi_{\varepsilon,x}=\varepsilon\Bigl(\sum_k x_kF_k+\beta(x)G\Bigr),\qquad
\beta(x)=\Bigl(\beta_0^2-\sum_k x_k^2\|F_k\|_2^2/\|G\|_2^2\Bigr)^{1/2}.
\]
Disjoint supports make $\|\Phi_{\varepsilon,x}\|_2$ independent of $x$, so the total incoming
mass does not depend on the target, and the tail $G$ pays for the change after every reading.

\begin{theorem}[Same past, open readings]\label{thm:same-past}
Fix $N\ge1$ and $0<T_1<\dots<T_N$. Subject to the external inputs below, there are $\mathcal E_*,c>0$ such that for every $0<\mathcal E<\mathcal E_*$ and every $y\in\mathbb R^N$ with $|y|<c\sqrt{\mathcal E}$ there is a spherically symmetric solution $(M,g,\phi,b,n)$ on $\mathbb R\times\mathbb R^3$ such that:
\begin{enumerate}
\item the entire pointed causal past $\Jm(b)$, with its field and launch tangent, is the same
Minkowski past with $\phi=0$ for all $y$;
\item the iterated incoming boundary mass $\lim_{V\to\infty}\lim_{u\to-\infty}m(u,V)$ equals $\mathcal E$;
\item the central future unit geodesic $\gamma$ with $\gamma(0)=b$, $\gamma'(0)=n$ exists for all
proper time, and $\phi(\gamma(T_j))-\phi(b)=y_j$ for each $j$.
\end{enumerate}
The selected solutions have smooth representatives that are future and past causally geodesically
complete, $C^0$-inextendible maximal globally hyperbolic developments.
\end{theorem}

\begin{proof}[Structure of the proof]
\emph{External inputs.} Theorem~1.1 and Definition~2.1 of \citet{LOY2016} give existence and
uniqueness on finite cones. Their (1.7)--(1.9), Proposition~3.18 and Section~4 give estimates
uniform in the cone vertex, and Theorem~1.7 with Section~6 gives causal completeness. Their local
well-posedness proof and part of the $u$-side higher estimates are omitted or sketched in the
source. The past half of their completeness theorem is stated in one sentence, and the extension to our compact family is our deduction from their uniform bounds. We retain all of these as inputs.

\emph{Flat wedge.} On $v<A$ the regular function $p=r\phi$ satisfies $p_{uv}=Vp$ with bounded $V$,
vanishing incoming data and a regular axis. An integral Grönwall argument gives $p=0$, hence
$m=0$ and the Minkowski geometry there. Along a future causal curve $v$ does not decrease, so
$\Jm(b)$ lies in the wedge.

\emph{Degree before the limit.} On each finite cone with vertex $u_n$ let $H_n(x)$ be the $N$
readings. Compactness and characteristic uniqueness make $H_n$ continuous. No uniform continuity in
$n$ and no continuous global scattering map is used. Explicit estimates give
$|\varepsilon^{-1}H_n(x)-x|\le C\varepsilon^2$ uniformly, so for small $\varepsilon$ the Brouwer degree
gives a root of $H_n(x)=\varepsilon z$ for each $n$.

\emph{Limits and mass.} Choose a convergent subsequence of roots and pass to a global reduced
solution by compactness, taking the cone limit before $u\to-\infty$. On a fixed finite interval
the exact mass flux identity then has integrand tending uniformly to $\Phi^2/4$, so
$\mathcal E=\|\Phi\|_2^2/4$. Choosing constants gives the quantifier order stated.

\emph{Regularity.} A finite-order reconstruction gives a $C^{1,1}$ metric and scalar in Cartesian
charts with unique geodesics. An all-order induction on the early exterior gives a smooth entire
slice. A polar-areal first-order system, with a constrained coefficient that gains a derivative,
propagates smoothness by localized $H^3$ persistence, weak--strong uniqueness and a tame continuation
estimate. Polar slices are Cauchy. Causal completeness with Theorem~1 of
\citet{GLS2018} gives $C^0$-inextendibility, and a field-preserving common extension
\citep{Ringstrom2009} gives maximality.
\end{proof}

Four qualifications matter. The constants depend on the schedule, so this is not one resolution for
all $N$; every finite sign pattern of readings occurs at $N$-dependent amplitudes. The fixed mass is
a total: the mass arriving by each reading time is not fixed. Incoming mass is not identified with
ADM mass. And the argument is target by target, so no scattering map is selected.

The readings differ because data from outside $\Jm(b)$ reach the center. That a past need not fix later readings is known in general \citep{Geroch1977,Manchak2008}. That is the expected mechanism, and the conclusion should not be oversold. For the linear wave equation on Minkowski
spacetime, $\phi(t,0)=\Phi(t)$ exactly, and every small central history is prescribable with the same
past. The class and its global regularity are in print \citep{Kehrberger2021}, and the method of
finite control by disjoint bumps and a topological theorem is established \citep{KehleUnger2022}. In
another category Carlotto and Schoen show that positive mass need not transmit data inward
\citep{CarlottoSchoen2016}. What Theorem~\ref{thm:same-past} adds is an exact statement of the
coupled form, with a mass and a complete maximal category held fixed.

A worry remains: perhaps these are one spacetime viewed from different marks.

\begin{proposition}\label{prop:uncountable}
Compare solutions with their central geodesic as a distinguished unparametrized curve $\Gamma$,
under time-orientation-preserving field isometries preserving $\Gamma$. If the central scalar trace is
$C^1$ and $N\ge2$, an open ball of readings cannot be realized by countably many isometry classes.
\end{proposition}

\begin{proof}
Such an isometry restricts to a translation of proper time on $\Gamma$. So every marked reading vector
from a copy of one solution lies on the curve $\beta\mapsto
(\phi(\gamma(\beta+T_j))-\phi(\gamma(\beta)))_j$, or its negative if sign reversal is allowed. That curve
is locally Lipschitz in $\beta$, so it has $N$-dimensional measure zero for $N\ge2$. A countable union
of null sets contains no ball.
\end{proof}

Bare metric nonisometry is a stronger claim and needs the center to be intrinsic.

\section{Whole-realm rigidity}\label{sec:rigidity}

Agreement on $\Jm(b)$ is not the only base a theory might hold fixed. Correia and Rosenkranz
distinguish a pointy from a bow-tie relativistic growing block \citep[Chap.~9]{CR2018}. In the notation of Section~\ref{sec:realms}, these are $E_b=\Jm(b)$ and $E_b=M\setminus(\Jp(b)\setminus\{b\})$.
The bow tie contains the elsewhere, and so contains exterior radiation that the causal past misses.
A past-agreement counterexample is not a counterexample to determination by that larger realm.

On a fixed background this is immediate. Put $W_b=\Ichr(\Ifut(b))\setminus\Jp(b)$. In general $W_b$ is smaller than the bow tie, which is why the sufficiency below is the stronger claim. If two solutions of a
law with the domain-of-dependence property agree on $W_b$, they agree on $\Jp(b)$: every past-inextendible
causal curve from a point of $\Jp(b)$ meets an early Cauchy slice at a point of $W_b$. This is
sufficiency, not a claim that all of $W_b$ is needed; exact dense data already determine continuous
fields, and the least region intercepting the relevant curves is not a minimal information base.

With gravity the comparison is between two spacetimes, and the theorem needs care.

\begin{lemma}\label{lem:cones}
In the class of Theorem~\ref{thm:same-past}, every two events share a chronological future. Writing
$W=M\setminus\Jp(b)$, every early level $\Sigma=\{t=t_0\}$ with $t_0<0$ is a Cauchy surface of $M$
contained in $W$, and is Cauchy in $W$.
\end{lemma}

\begin{proof}
In the global Cartesian coordinates of the $C^{1,1}$ reconstruction, whose causal curves are those of the smooth representative, the metric has $g_{tt}=-F$, no time--space terms, and spatial
eigenvalues $F$ and $\mathcal A$ with $2/9\le F\le8/3$ and $4/9\le\mathcal A\le1$. So $(1,v)$ is timelike for
$|v|\le1/4$, and straight paths to $(T,0)$ with $T$ large give a common chronological future. For the
second claim, $t$ increases along future causal curves and $t(b)=0$, so $\Sigma\subset W$. The set $W$ is an
open past set, so a past causal curve starting in it stays in it, and an inextendible curve of $W$
cannot end at finite $t<t_0$, since the only ambient boundary of $W$ lies in $t\ge0$, nor escape to spatial infinity, by the cone bound.
\end{proof}

Since every two events share a chronological future, $W_b=W$ in this class, and $W$ is the bow tie minus $b$.

\begin{theorem}[Whole-realm rigidity]\label{thm:rigidity}
Let $M_1,M_2$ be smooth maximal developments in this class with marks $b_1,b_2$. Suppose
$h:W_1\to W_2$ is an onto time-orientation-preserving isometry of metric and scalar. Then there is an onto smooth field-preserving isometry $\Psi:M_1\to M_2$ extending $h$, with $\Psi(b_1)=b_2$. If $h$ also matches the launch germ and tangent, all
corresponding future proper-time readings agree.
\end{theorem}

\begin{proof}
Let $\Sigma_1$ be as in Lemma~\ref{lem:cones} and $\Sigma_2=h(\Sigma_1)$. Then $\Sigma_2$ is Cauchy in $W_2$ and
acausal in $M_2$, since $W_2$ is causally convex. Every ambient inextendible causal curve has a nonempty past
tail in $W_2$, and its maximal connected portion there is inextendible in $W_2$, so it meets $\Sigma_2$. Hence
$\Sigma_2$ is Cauchy in $M_2$. The isometry matches the induced metric, extrinsic curvature, scalar and its normal
derivative on these surfaces. Both spacetimes are maximal developments of those data, so uniqueness \citep{CBG1969,Ringstrom2009}, in the scalar form of \citet{RingstromNote}, gives an onto smooth field-preserving isometry $\Psi$ with $\Psi|_{\Sigma_1}=h|_{\Sigma_1}$. $\Psi$ and $h$ have the same first jet there, so they agree on the connected set $W_1$, and $\Psi$ carries $\Jp(b_1)=M_1\setminus W_1$ onto $\Jp(b_2)$. Thus $\Jp(\Psi(b_1))=\Jp(b_2)$, and $\Psi(b_1)=b_2$ since $M_2$ is causal. The last claim follows from uniqueness of geodesics.
\end{proof}

In one physical class, with one deterministic law:
\begin{center}
\begin{tabularx}{\textwidth}{@{}Yl@{}}
\toprule
Held fixed & Future central readings\\
\midrule
Whole pointed causal past, and total incoming mass & Open ball\\
Whole open realm $W$, with marked germ & Fixed, and so is the whole development\\
\bottomrule
\end{tabularx}
\end{center}
The word ``present'' cannot be substituted for both bases. Neither result says which, if either, is
what exists from $b$, and the second does not make the determined future exist. It does show that a theorist of the bow-tie realm has a future fixed by what exists from $b$, while a theorist of the pointy realm in the same class does not. That holds for the bow tie generally, since it contains a Cauchy slice; the theorem's added content is the transport between two spacetimes. Openness relative to a realm depends on the realm chosen and the law, and the choice of a realm is not forced by physics.

\section{Realm structure}\label{sec:realms}

Which realms are available? Write $E_b$ for a point-level assignment over a fixed causal frame.
This is a metalinguistic device, and implies no claim that the frame is concrete. Two conditions do
much of the work: $b\in E_b$, and persistence along the causal order, $b\le c\Rightarrow E_b\subseteq E_c$.

\begin{proposition}\label{prop:sp}
Self-inclusion and causal persistence imply $\Jm(b)\subseteq E_b$. With future exclusion
$E_b\cap\Jp(b)=\{b\}$ and no elsewhere point this gives $E_b=\Jm(b)$; with every elsewhere point it gives
the bow tie.
\end{proposition}

\begin{proof}
If $x\le b$ then $x\in E_x\subseteq E_b$. The rest is the partition of $M$ into past, future and elsewhere.
\end{proof}

No field equation, Lorentz invariance or transitivity enters. The persistence condition matters:
with persistence only along timelike order, $\Ichr(b)$ is available too. Correia and Rosenkranz's
clauses use the full causal order \citep[pp.~144, 148--149]{CR2018}. A presentist denies persistence.
Choosing between pointy and bow-tie needs further premises. Requiring existence realms to be closed
would favour the pointy profile, but metric regularity does not supply that; it is a bridge.

The bow tie is also not always what a natural criterion gives. Let $D(b)$ consist of $b$ and the points
outside its strict causal future that share a strict causal successor with it. On connected causally simple spacetimes this is persistent exactly when every pair of events has a common strict causal future. In the class of
Theorem~\ref{thm:same-past}, Lemma~\ref{lem:cones} gives $D(b)$ equal to the bow tie. Elsewhere it need
not be: on a complete conformal cylinder of duration $L$ over a complete spatial manifold $S$, the least
persistent enlargement of $D$ equals the bow tie at every standpoint exactly when
$\operatorname{diam}(S)\le L$. Whether a horizon obstructs the criterion depends on the class.

\partheading{Part IV. Existence and truth}

\section{Stages, ended worlds and maximality}\label{sec:maximal}

A truncated spacetime can represent two different things: what exists so far, on a growing view, or a
whole world that stops. Conflating them makes an openness claim an artefact of representation.

Let $(M,g,\phi)$ be a smooth member of the class above with global center-normalized time $T$, and fix
$T_0<T_*$. Then $M_<=\{T<T_*\}$ is an open smooth solution with the same complete Cauchy data on $\{T=T_0\}$,
whose constant-$T$ slices remain Cauchy, and for $T_*$ after $b$ the pointed past of $b$ is unchanged. The
field equations fail nowhere on it. So three principles have different strengths:
\begin{description}
\item[Local lawfulness:] every included point satisfies the equations. Both $M$ and $M_<$ do.
\item[Continuation is possible:] a regular end with no boundary obstruction admits a lawful extension.
\item[Maximal-world admissibility:] complete histories representing physically possible worlds are maximal
developments, unless a physical end rule is stated.
\end{description}
Only the third excludes $M_<$ as a complete world, and it is not a consequence of the first. It is a
defensible methodological choice for comparisons that claim to hold a theory fixed: choosing a shorter
domain adds a contingent restriction with no boundary mechanism. It is not a refutation of brute endings.

Accepting maximality concedes no ontology. A growing theorist can use $M$ to represent everything that
does, did or will happen while assigning a proper inventory to each standpoint. So growth with a physically
determined future is consistent at this level, which refutes any inference from growth alone to
indeterminism. A completed mathematical representation and an ontological completed history are different
commitments. The causal-set case makes the same point sharply: the prescribed complex percolation values
extend to a countably additive measure on the labelled doubly-principal cylinder algebra exactly when
$p\in[0,1]$ \citep{ProjectPercolation}. A failure of completion in a growth theory is a mathematical fact about a
measure. It neither refutes becoming nor, where it succeeds, establishes a block.

\section{Truth}\label{sec:truth}

Fix a claim $q$ at $b$ about a reading at a later event. Four claims should be separated. \emph{Outcome
existence}: in each admitted complete realization the target occurs with a definite value.
\emph{Settlement}: every realization admitted by the chosen base gives $q$ the same value. \emph{Truth now}:
$q$ has an absolute truth value at $b$. \emph{Grounding}: an appropriate fact explains that truth, perhaps
under a later temporal operator.

Take two realizations sharing a stage at $b$, with the target absent from the inventory at $b$ and included
when it occurs, and with the value differing between them. On policy L, evaluated in a designated realization,
$q$ is bivalent at $b$, and its truth can be assigned a later ground. This is the restricted form of the later-grounded view \citep{CR2018}, which \citet{Wawer2021} challenges. On policy S, $q$ is evaluated over both realizations and is neither
supertrue nor superfalse at $b$. Both policies have eventual occurrence, a definite eventual value and the
same change of inventory. They disagree about truth now. So eventual resolution with growth does not entail
truth now, and it does not exclude it. This is a restricted counterinterpretation, not a model of any named
theory or a refutation of bivalence, and the logical point belongs to an established debate.

The physics fixes the comparison and not the policy. Theorem~\ref{thm:same-past} supplies realizations sharing a
pointed past. Which of them count as admitted, and whether one is designated, is the policy. Designating an
actual realization is not derived from the past of $b$.

\section{The direction of time}\label{sec:direction}

Nothing so far says why the past should be the base. Write $F^-_{\mathcal L}(b)=F_{\mathcal L}(\Is(b))$ and $F^+_{\mathcal L}(b)=F_{\mathcal L}(\Isp(b))$, and let $\Det^+(b)$ be the set of events all of whose future-inextendible causal curves meet $\Isp(b)$.

\begin{proposition}[Mirror]\label{prop:mirror}
Let $\theta$ be an isometry of $M$ reversing time orientation, and let $\mathcal L$ be invariant under pullback by $\theta$. Then $\theta(F^-_{\mathcal L}(b))=F^+_{\mathcal L}(\theta(b))$ and $\theta(\Det(b))=\Det^+(\theta(b))$. Without such an isometry, every result of Parts~II and~III that is stated for a time-reversal invariant law holds for the same spacetime with its time orientation reversed.
\end{proposition}

\begin{proof}
A difference of solutions vanishing on $\Is(b)$ pulls back to one vanishing on $\Isp(\theta(b))$, and $\theta$ exchanges past- and future-inextendible causal curves. The last claim holds because those results use the orientation only as a label.
\end{proof}

The proposition is elementary, and its consequences are not idle. In Minkowski spacetime the conformal scalar is fixed from $b$'s strict future exactly on $\Jp(b)$. In the closed dust model the reflection $\eta\mapsto2\pi-\eta$ is such an isometry, so $\Det^+(b)=M$ for $\eta_b<\pi$. The statement that the guarantee changes with cosmic epoch is true, and it has an exact mirror. The Einstein--scalar system is time-reversal invariant, so Theorem~\ref{thm:same-past} has a twin in which a causal future and total outgoing mass leave earlier central readings free. Theorem~\ref{thm:main} carries its orientation in (CB); reverse it and the same proof gives the future lobe.

So exact fixing has no arrow of its own. Under equal bases each side is as fixed or unfixed from the other. ``The past is fixed and the future open'' cannot be read off the fixing relation, and the asymmetry needs another source. We know of three candidates, and the paper is neutral between them. Each gives (CB) a different meaning.

\emph{Thermodynamic.} The dynamics is reversal invariant, and statistics over a present macrocondition make entropy higher on both temporal sides. So Albert adds a law-like posit, the Past-Hypothesis, that the world began in a particular low-entropy macrocondition \citep[p.~96]{Albert2000}. Prediction uses the dynamics and the statistical rule, while most of what we know of the past rests on records, and that asymmetry is meant to follow from the posit \citep[pp.~113--116]{Albert2000}. On this view the past is a natural base because it is what we have records of, and that is explained by a boundary condition and not by $\mathcal L$.

\emph{Perspectival.} Price warns against accepting an argument in one temporal direction while rejecting its mirror \citep[chs.~1--4]{Price1996}. He separates the asymmetries of dependence, causation and agency \citep[pp.~132--133]{Price1996} and defends a perspectival view of causal asymmetry. He presupposes that the present and flow are perspectival, so he is a party to the dispute. Proposition~\ref{prop:mirror} is the fixing instance of his warning. His treatment of cosmology also bears on Part~II: if a low-entropy boundary needs explaining at one end, a time-symmetric Gold universe cannot be dismissed by arguments that assume the arrow \citep[p.~81ff]{Price1996}.

\emph{Causal.} Tooley takes causal direction as basic and bases temporal direction on it \citep{Tooley1997}. A cause is actual as of its effect's time and not conversely \citep{Tooley2010}. On this view the asymmetry of openness would be ontological. He also defends absolute simultaneity, which is the clock-equipped horn of Section~\ref{sec:becoming-rival}.

One question stays open. If $\mathcal L$ includes the Past-Hypothesis it is not reversal invariant, and Proposition~\ref{prop:mirror} fails. We expect exact fixing to change little, since a macrocondition does not fix exact values, and the asymmetry to appear in conditional probabilities over completions. If so, chance given a base belongs beside exact fixing, and much ordinary talk of openness may mean that. We have not proved it.

\section{The upper bound and the lower bound}\label{sec:bounds}

\subsection{What each side owes}

Let $R_b$ be what a restricted view counts as concrete from $b$, and $H$ the represented history. Constitution gives $R\subseteq E_b$ for an activity region $R$, which may include events later than $b$ (Section~\ref{sec:process}). Persistence gives $\Jm(b)\subseteq E_b$ (Proposition~\ref{prop:sp}). A restricted view adds $R_b\subseteq E_b$. These are lower bounds, and every argument from physical organization that we have examined stops at one. Theorem~\ref{thm:main} gives an upper bound only by amplifying an assumed failure of universality, as discussed below. A lower bound cannot settle the matter, because $R_b\subseteq E_b$ is compatible with both $E_b=R_b$ and $E_b=H$.

\begin{center}
\begin{tabularx}{\textwidth}{@{}lYYY@{}}
\toprule
 & Pointy restricted & Slice or bow-tie restricted & Event realist\\
\midrule
$E_b\subseteq H$ & Holds for every completion through $b$ & Only for completions agreeing on the slice & Holds of one $H$\\
Selection owed & None & For that part & Which completion, where they differ\\
Extra claim & $E_b\subseteq R_b$, negative, differs by $b$ & The same & $H\subseteq E_b$, positive, the same for every $b$\\
Must say & Whether more will come, or why that is unsettled & The same & Why $H$ and not a truncation\\
Supported by $R_b\subseteq E_b$? & No & No & No\\
\bottomrule
\end{tabularx}
\end{center}

So the two debts are not of one kind. Where $H$ is unique in a class all parties accept, the event realist's extra
claim is a stronger lower bound. Where completions are not unique, the event realist's $E_b\subseteq H$ is a selection
that a pointy theorist does not owe. Three views fall outside the table: on McCall's view as reported by \citet{Earman2008} $E_b$ lies
in no single history, fragmentalists have no one $H$, and for a theorist whose $H$ has no referent beyond $R_b$,
``nothing concrete lies outside $H$'' is the thesis itself. The positive claim carries more events and the negative
one a changing limit. We know of no neutral way to price that exchange.

\subsection{Totality}

Is ``that's all'' an instance of a known problem? Armstrong's totality fact relates a mereological whole to a property when the whole comprises all things with the property, and serves as truthmaker for negative and general truths \citep[\S6.2]{Armstrong2004}. The property may be abundant and time-indexed, and he grants that the relation sets a limit \citep[pp.~72--73]{Armstrong2004}. Molnar argues it is a negative fact under another
name \citep{Molnar2000}. The rival accounts differ only about the truthmaker: absences, the world with its contents
essential \citep{Cameron2008}, the whole positive fact \citep{CheynePigden2006}, or no truthmaker, on optimalist views with Bigelow's supervenience of truth on being (both as reported in \citealp{MacBride2024}). All start from an accepted negative truth and ask what grounds it. None
is a reason to accept it. Filled as $T(R_b,\text{concrete from }b)$, the slot restates the claim; if the property is
membership in $R_b$, it is trivial. Forrest's use of general facts in a growing block \citep{Forrest2006} appears to do other work, but we have not read it beyond its abstract and one review, and it could change the first row of the table below.

The neighbouring literature has the core points. Earman says a growing-block theorist cannot deny the future absolutely
without fixing one truncated model and giving up growth, so the denial is perspectival, and a block theorist can
say the perspectival version as a triviality \citep[pp.~144--145]{Earman2008}. He separates this rigid denial from a
shifting one, that far enough beyond now nothing exists, which holds at every stage and fails in block models
\citep[pp.~144, 159]{Earman2008}. His ``absolute'' is not Fine's: the standard realist holds a tensed denial absolutely
and keeps growth through what will be, and Deng rejects the step from that view to one fixed block
\citep[\S4]{Deng2017}. Deng also objects to Tooley that what is actual as of now includes that those facts are all the
tenseless facts, and that this is tensed \citep[\S6]{Deng2017}. None of these sources discusses the inference from
constitution to inventory, so they do not state the gap above. No priority audit has been run.

\subsection{Persistence and the rigid claim}

Persisting truths are closed under entailment, so a claim that fails to persist has no persisting entailers.
``Everything lies in $R_b$'', with $R_b$ rigid, fails to persist wherever the inventory changes, on a growing block or on
presentism. That is the reason behind Earman's rejection of a rigid denial. It does not touch the shifting form
``everything is past or present''. The rigid instance follows from the shifting form plus self-location, so the only
non-persisting ingredient is locational, and the debt moves to the shifting form with an unrestricted quantifier. The
observation covers positive tenseless and past-tensed premises. It leaves open past-tensed conclusions and non-deductive
support. It does not cover totality-form constitution claims or a presentist's present-tensed premises. It is close to the familiar point that tenseless premises do not yield tensed conclusions, and it says
nothing against the event realist.

\begin{proposition}\label{prop:tetrad}
A growing-block theory cannot hold truthmaker maximalism, temporal necessitation, permanence of what exists, and a rigid
``that's all'' truth at $b$ which growth falsifies.
\end{proposition}

\begin{proof}
This is Loss's argument \citep[\S II]{Loss2022}, with growth in place of A-change. If $p$ is true at $b$ it has a
truthmaker, which by permanence exists at a later stage, and by necessitation $p$ holds there, contrary to growth.
\end{proof}

Given the shifting form, the rigid claim is equivalent to ``$b$ is present'', which is Loss's own case, and his \S V treats
presentness on the block as being at its edge. Presentism lacks permanence and the block view lacks change, so neither
is touched. The exits are these.

\begin{center}
\begin{tabularx}{\textwidth}{@{}lY@{}}
\toprule
Member given up & Cost\\
\midrule
Permanence & A fact ceases at each step, against Broad's text and Loss's (GB2)\\
Temporal necessitation & Truthmakers gain proper parts, the block is the truthmaker and mereological universalism goes
\citep[\S\S V--VI]{Loss2022}. Molnar's rejection of contingent truthmaking \citeyearpar[\S IX]{Molnar2000} targets a primitive
relation; whether it reaches this version is unargued\\
Maximalism & Growth stays additive, and the upper bound is no fact at the fundamental level. Cameron and Cheyne and Pigden argue
against the restriction\\
Truth-value change & Index the proposition, as Tooley does \citeyearpar{Tooley1997}. Actuality as of a time is primitive, and Deng's
objection applies\\
\bottomrule
\end{tabularx}
\end{center}

Separately, each view takes a stance on Fine's four postulates \citep{Fine2005}, as Deng sets them out. Those stances are not exits
from this table: relativizing every clause to standpoints reproduces the contradiction at the later standpoint, and a standard realist
faces it with tense operators. Nothing above needs a foliation, since a fact ceases at $s$ if it holds at $s$ and fails at every later
stage, which needs only a partial order. What relativity removes is a privileged chain of stages, and without one, parity pushes a view
towards relativism. That would hold in Newtonian spacetime too.

\subsection{Stated candidates}

The following table sorts the stated upper-bound candidates we know. It covers these candidates only; it is not a theorem about
all such principles.

\begin{center}
\begin{tabularx}{\textwidth}{@{}YlY@{}}
\toprule
Candidate & Index & Outcome\\
\midrule
Armstrong totality & None & Restates the claim\\
Maximal-world admissibility & None & Bounds a complete history, not $E_b$\\
Theorem~\ref{thm:main} via the reality bridge & Defined or primitive & Conclusion shared if defined; non-universality rejected if primitive\\
Interaction closure \citep{Brading2015} & None & Inclusive as stated; with its converse it falls at $H$ if connected\\
Causal criterion, tenseless \citep{Colyvan1998} & None & Later events pass\\
Causal criterion, indexed & Defined or primitive & The same bound as Theorem~\ref{thm:main}, given the bridge and non-universality\\
Tooley: an effect is not actual as of its cause's time & Primitive & Rival rejects. Support comes from an analysis of causation, and the direction it needs is disputed (Section~\ref{sec:direction}). Paired with absolute simultaneity\\
No fact selects a completion, the reply Earman reports \citep[pp.~144, 154, 159]{Earman2008} & Primitive & Rival rejects the bridge. Silent in deterministic classes\\
Least truthmaker base, given settled truth & Primitive & Abductive. Rival rejects the truth policy. Our construction\\
\bottomrule
\end{tabularx}
\end{center}

Theorem~\ref{thm:main} is a real upper-bound derivation, of a conditional kind. It amplifies an assumed failure of universality into a
sharp bound and does not supply the first negative. The bound is on what has happened by $b$, and an eternalist can accept that.
Existence enters only at a bridge, which reads having become as being real.

The causal criterion deserves a remark, because it is the closest thing to a sourced upper bound. The Eleatic principle admits only
what has causal power. Read tenselessly it passes every later event with causes or effects, so it gives no restriction at the future.
Only a standpoint index does: ``concrete from $b$ only if it can affect $b$''. That is the converse of Clifton and Hogarth's maxim,
and the maxim's rationale, that an unreal event could not help bring $b$ about, gives no support to the converse. Colyvan notes that the
epistemic motivation for the principle gives no reason to believe in future objects, and counts a similar exclusion of distant regions
against that motivation \citep{Colyvan1998}.

\begin{quote}
\textbf{The diagnostic.} An index on ``concrete from $b$'' can be defined or primitive. Defined as a causal or temporal restriction, the
claim is a triviality every party accepts. Primitive, it is a thesis the event realist rejects, indexed or not.
\end{quote}
So ``indexed'' sorts nothing. Correia and Rosenkranz's clause (P2R) and Tooley's actuality as of a time are indexed and rejected. Static
views give no sense to indexed actuality \citep{Tooley2010}. For Fine's relativist, what is said from $b$ still carries the index, so
``absolute'' is not available as a contrast. The missing argument lies between the defined and primitive readings. Three of the entries above, Tooley's,
the no-fact principle and the least-base argument, offer support of a different kind, and they fall outside this sorting.

\section{Constitution and process}\label{sec:process}

One physical route to a restricted inventory runs through systems. If a genuine unity is constituted by an activity
extending through a region $R$ \citep{Brading2015}, perhaps $R$ is what exists for that system. The inference fails at one identifiable step.
Grant that the unity's realizing occurrences must be concrete. Compare $E_b=R$ with $E_b=H$. Both satisfy $R\subseteq E_b$, so an argument
whose only inventory requirement is that constituting occurrences are concrete cannot distinguish them. It supplies a lower bound. A process can
have exactly its own parts while other events are also concrete. So the route from constitution to an exhaustive inventory is closed as a deductive route unless a further exhaustiveness premise is supplied, and adding examples will not supply it. The revision that makes the activity the object's parts earns a process mereology, and
the rival can accept that.

An explanatory route remains: does a growing inventory explain process unity better? Take a feedback-controlled filling device with a sensor, a
stored target and a valve, and compare a timed open-loop replay that closes at the same moment. For one supply rate the two have identical
liquid height and flow histories. Under a changed rate the feedback device reaches its target later, and the replay stops short. So organized causal
roles explain something a list of outputs does not, and a process is not a rigid aggregate of its actual events. But an event realist who
identifies the operation with the causally organized execution, with an explicit participation rule, keeps all of that, and so does a growing
theorist. The argument for an exclusive advantage assumed that event realism must identify a process with a rigid aggregate, and it need not
\citep{Stout1997,Steward2013}. This is a comparison for one explanandum, not a verdict of parity, and whether a further form or grounding fact is
needed remains open.

\section{Which inferences fail}\label{sec:revise}

The results license no ontology, but they do bear on particular inferences. None of what follows is claimed as new.

\begin{description}
\item[Metric-only shared presents.] Corollary~\ref{cor:nopresent}. A transitive relation containing the chronological past and defined from the metric alone has no shared present if it is local. Its defenders must name the premise they reject: locality on $\GH$, transitivity, the off-shell premise, or the category.
\item[The causal past as a base for every rival.] Theorem~\ref{thm:rigidity}. A comparison must match the rival's realm. A past-agreement counterexample does not reach a bow-tie view, whose realm contains a Cauchy slice.
\item[Causal limits on prediction read as openness.] Theorem~\ref{thm:map}. What is guaranteed does not bound what a law fixes, as \citet{Manchak2008} and \citet{Earman1995} partly anticipate.
\item[``Past fixed, future open'' as evidence for growth.] Proposition~\ref{prop:mirror}. The slogan cannot be read off the fixing relation. The asymmetry of records has a candidate explanation that uses no inventory distinction, so it is not evidence for growth until that explanation is shown to fail. The growing theorist can reply that it does fail. This removes one support and refutes nothing.
\item[Growing-block maximalists.] Proposition~\ref{prop:tetrad}. This is Loss's argument, and it fires through the rigid claim.
\item[Unreality to law-relative openness.] Theorem~\ref{thm:map}, on the reading of having become as being real. The distinction is old; the models are exact. Openness of existence is untouched, since it concerns the modal class (Section~\ref{sec:maximal}).
\item[Past-relative nonfixing to growth.] Theorem~\ref{thm:same-past} holds for an eternalist, since each model is one whole spacetime with its own exterior data. It is no evidence against indeterminist growth, and we know of no one who argues this way from it.
\item[Reality to truth.] \citet[pp.~243--246]{Putnam1967} argues that future contingents already have truth values. That step needs a policy connecting outcomes to earlier truth (Section~\ref{sec:truth}). A block theorist can adopt the policy at no cost to the inventory, so the reason is to state it.
\item[Constitution to inventory.] Section~\ref{sec:process}, against the argument as we reconstructed it.
\item[Event realism as a default.] Where completions are not unique, the selection of $H$ is a debt. A selection principle preferring fewer primitives \citep{ProjectRole} needs a defence of its own.
\end{description}

One reading of \citet[pp.~148--149]{Stein1991} is also excluded. If ``definite'' and ``unsettled'' meant fixed by law, the relation would not be $\Jm$, since Theorem~\ref{thm:map} has $F=M$ in several rows with $B$ inside the causal past. Nothing we have read shows that Stein intends that reading, and his theorem does not need it.

The positions compare as follows. Each row is a selected profile, not a reconstruction of a whole theory.

\begin{center}
\begin{tabularx}{\textwidth}{@{}lYYY@{}}
\toprule
View & Concrete from $b$ & Premise denied or added & Future fixed from what exists?\\
\midrule
Presentism & Present only, however spatially read & Denies persistence; owes an upper bound & Depends on the spatial reading and the law: as the clock row for a Cauchy slice; past-grounding owed\\
Pointy growing block & $\Jm(b)$ & Owes an upper bound & Open in Theorem~\ref{thm:same-past}; fixed in closed dust after maximum expansion\\
Bow-tie growing block & Past and elsewhere & Rejects transitivity of existence-at; owes an upper bound & Fixed wherever the law has unique evolution, since the realm contains a Cauchy slice\\
Standpoint relativism & $\Jm(b)$, relative to $b$ & Rejects Absolutism & As pointy; the geometry is agreement among standpoints\\
Clock-equipped view & $\{T\le T(b)\}$, with a level of $T$ as present & Changes the category & Fixed if the law has unique Cauchy evolution\\
Event realism & $H$ & Owes a lower bound, and selection where completions are not unique & Yes, trivially, and that says nothing about law\\
\bottomrule
\end{tabularx}
\end{center}

\section{What would change these conclusions}\label{sec:defeat}

A gluing counterexample, or a category mismatch in Part I. A proof that on-shell energy conditions change the classification; the
pocket problem of Section~\ref{sec:energy} is the route. A failure of the external inputs of Theorem~\ref{thm:same-past}. A prior theorem with the same
hypotheses, which would change credit and not content. In the philosophy: an independently defended tensed premise yielding the upper bound; an
explanandum for which a growing inventory does exclusive work; or an argument that the defined/primitive contrast is unavailable, as a
deflationist about the debate may hold \citep{Savitt2006,Dorato2006}. A successful deflation would remove the existence question and leave Parts I to III standing.

\section{Conclusion}\label{sec:conclusion}

The future can be unreal and fixed, real and unfixed from any standpoint, fixed at its end before its middle, and open or closed in one physical class
according to what is held fixed. General relativity with a locality premise says sharply what can have become. Laws and comparisons say what is fixed, and in the models treated we can say it exactly. Neither says what exists, and neither says what is true. What the past can fix turns out to be a question physics can answer, and it is a
different question from the one it is usually taken to answer.

\paragraph{Note on method.} This manuscript was developed with substantial assistance from large language models (Anthropic Claude, OpenAI Codex and GLM models), including
literature search, proof drafting and cross-model proof checks that found and repaired errors. Review by models is internal review. It is not expert or
journal review, and agreement between models is not evidence of correctness. Responsibility for errors is the author's.

\appendix

\section{Proofs for Part I}\label{sec:proofs}

\subsection{Theorem \ref{thm:class}, Theorem \ref{thm:weyl}, Proposition \ref{prop:reg}}

We use: $\Ifut_{\mathrm s}$ is open in $M\times M$ (MS Prop.~2.16);
$\overline{\Jp}=\overline{\Ifut_{\mathrm s}}$ and $\operatorname{Int}\Jp=
\Ifut_{\mathrm s}$ (MS Prop.~2.17), and likewise for the pasts; if $p\ll r\le q$
or $p\le r\ll q$ then $p\ll q$, and if $r\in E^+(p)$, $q\in E^+(r)$ then
$q\in E^+(p)$ or $p\ll q$ (MS Prop.~2.31); two events joined by a causal
curve that is not a maximizing lightlike curve are joined by a timelike curve,
and a maximizing lightlike curve is a lightlike geodesic (MS Thm.~2.30(i),
(ii)); so a causal curve between horismotically related events is a lightlike
geodesic.

\begin{proof}[Proof of Theorem \ref{thm:class}]
By Theorem~\ref{thm:main}(a), $B=\Ichr\cup N$. Timelike and causal curves
between points of a causally convex set stay in it, so $\Ichr$, $\Jm$ and $E^-$
are locally covariant and $N$ inherits (LC); (NC) is transitivity restricted
to horismotic pairs. Conversely, $\Ichr\cup N$ is reflexive and satisfies (CB);
for transitivity, a composition with a chronological step is chronological by
MS 2.31, and a composition of two $N$-steps is chronological or horismotic, in
which case (NC) applies. The last claim is Theorem~\ref{thm:flat}.
\end{proof}

\begin{proof}[Proof of Theorem \ref{thm:weyl}]
(NC): if $x\to y\to b$ along lightlike segments $\gamma_1,\gamma_2$ and
$x\in E^-(b)$, then $\gamma_1\gamma_2$ is a causal curve between horismotic
points, hence an unbroken lightlike pregeodesic (MS Thm.~2.30), and it
contains the Weyl point of $\gamma_1$. (LC): lightlike geodesics between points
of a causally convex image stay in it, and isometries preserve the Weyl
tensor. On $\Mink$ the Weyl tensor vanishes. If it vanishes nowhere, every
horismotic pair is in $N$. If $M$ is flat on an open causally convex $V$ and
the Weyl tensor is nonzero at some point $z$, horismotic pairs inside $V$ are
not in $N$ (every causal curve between them stays in $V$), while a horismotic
pair joined by a lightlike segment through a neighbourhood of $z$ is. The Weyl
tensor, as a $(1,3)$ tensor, is invariant under conformal rescaling (standard),
and unparametrized lightlike pregeodesics are conformally invariant (MS
\S2.6). The last case is not vacuous: the spacetime $M'$ of
Section~\ref{sec:proofs}, built from a $g$ with nonzero Weyl tensor at $x$, is
flat on the open causally convex set $\{t<-3r\}$.
\end{proof}

\begin{proof}[Proof of Proposition \ref{prop:reg}]
If $B_M$ is closed and contains $\Ichr$, it contains
$\overline{\Ichr}=\overline{\Ichr_{\mathrm s}}=\overline{\Jm}\supseteq\Jm$, so
$B=\Jm$ by Theorem~\ref{thm:main}(a). If $B_M\setminus\Delta$ is open and
contained in $\Jm$, it lies in $\operatorname{Int}\Jm=\Ichr_{\mathrm s}$, so
$B_M\subseteq\Ichr_{\mathrm s}\cup\Delta=\Ichr$, and (CB) gives $B=\Ichr$.
\end{proof}

\begin{proof}[Proof of Corollary \ref{cor:worldline}]
By worldline becoming, $b\in B_{M,\lambda'}(b)$ for every $\lambda'\ni b$; with
$c=b$, worldline transitivity gives $B_{M,\lambda}(b)\subseteq B_{M,\lambda'}(b)$,
and by symmetry they are equal. Call the common set $B_M(b)$. Transitivity: if
$a\in B_M(b)$ and $b\in B_M(c)$, choose $\lambda\ni b$, $\lambda'\ni c$ and apply
worldline transitivity. (CB): if $a\ll b$, a timelike curve from $a$ to $b$
extends to an inextendible timelike curve $\lambda$, and $a\in\Ichr(b)\cap\lambda
\subseteq B_M(b)$. (LC): for $\psi:N\to M$ take $\lambda_N\ni b$ and an
inextendible extension $\lambda$ of $\psi\circ\lambda_N$; then $a\in B_N(b)\iff
a\in B_{N,\lambda_N}(b)\iff\psi(a)\in B_{M,\lambda}(\psi(b))\iff\psi(a)\in
B_M(\psi(b))$.
\end{proof}

\subsection{Theorem \ref{thm:main}}

The idea: a forbidden pair in a curved spacetime is moved, by local
covariance, into an auxiliary spacetime that agrees with the original near
the pair and is flat elsewhere; chronological becoming and transitivity then
produce a forbidden pair in a flat region, contradicting
Theorem~\ref{thm:flat}.

\begin{lemma}[Future pairs reduce to spacelike pairs]\label{lem:R}
If $x\in\Jp(b)\setminus\{b\}$, there is $e\ll x$ spacelike to $b$.
\end{lemma}

\begin{proof}
Let $U$ be a convex normal neighbourhood of $b$ and $V\subseteq U$ an open
neighbourhood of $b$ that is causally convex in $M$ ($M$ is strongly causal;
MS Remark~2.12 and the construction following it). Causal relations in $U$,
regarded as a spacetime, are read off $\exp^{-1}$ (MS Prop.~2.10), and for
points of $V$ they agree with those of $M$ (MS Remark~2.12). A future causal
curve from $b$ to $x$ is nonconstant ($M$ is causal), so it contains some
$b_1\in V\setminus\{b\}$; then $b\le b_1\le x$ and $u=\exp_b^{-1}(b_1)$ is
a nonzero future causal vector. Choose a spacelike $a\in T_bM$ with
$g(u,a)>0$ and put $q_\varepsilon=\exp_b(\varepsilon a)$. Synge's world
function $\sigma$ is smooth on $U\times U$ with $\nabla_q\sigma(q,b_1)|_{q=b}=-u$,
so $2\sigma(q_\varepsilon,b_1)=g(u,u)-2\varepsilon g(u,a)+O(\varepsilon^2)<0$ for
small $\varepsilon>0$. The vector $\exp^{-1}_{q_\varepsilon}(b_1)$ depends
continuously on $\varepsilon$ and tends to $u\neq0$, so for small $\varepsilon$
it is future-directed timelike: $q_\varepsilon\ll b_1$. Also
$2\sigma(q_\varepsilon,b)=\varepsilon^2g(a,a)>0$, so $q_\varepsilon$ is spacelike
to $b$. By MS Prop.~2.31(i), $e=q_\varepsilon\ll x$.
\end{proof}

The same local argument shows that every point has spacelike-separated
points arbitrarily close to it.

\begin{lemma}[Spacelike chains]\label{lem:E}
In connected globally hyperbolic $M$, any two points are joined by a finite
chain of consecutively spacelike-separated points.
\end{lemma}

\begin{proof}
Globally hyperbolic spacetimes (MS Def.~3.70) are causally simple (MS
Prop.~3.71), so $\Jp(z)$ and $\Jm(z)$ are closed and the set of points
spacelike to $z$ is open. Fix $p$ and let $S$ be the set of points reachable
from $p$. $S$ is open: if $q\in S$, pick $z$ spacelike to $q$; a
neighbourhood of $q$ is spacelike to $z$ and reachable via $q,z$.
$M\setminus S$ is open: if $q\notin S$ and $c$ is spacelike to $q$, a
neighbourhood $W$ of $q$ is spacelike to $c$; were some $q''\in W$ in $S$, then
$c\in S$ and $q\in S$. $M$ is connected.
\end{proof}

The example $M=\mathbb R\times S^3$, $b=(0,p)$, $x=(3\pi,p)$ shows that a single
point spacelike to two causally related points need not exist; the chain is
needed.

\paragraph{Set-up.} Fix spacelike-separated $x,b\in M$. Since $J^+\subseteq
M\times M$ is closed (MS Prop.~3.68(iii), 3.71), there are small closed
spacelike embedded $3$-discs $D_x\ni x$, $D_b\ni b$, each acausal, with no causal
curve between them; so $D_x\cup D_b$ is a compact acausal spacelike
submanifold with boundary, and Bernal and Sánchez give a smooth spacelike
Cauchy hypersurface $S\supseteq D_x\cup D_b$ \citep[Thm.~1.1, Rem.~4.14]{BS2006}.
Choose $\delta_0>0$ and a smooth embedded curve $c:[-\delta_0,L+\delta_0]\to S$,
unit speed for the induced metric $h$, with $c(0)=x$, $c(L)=b$.

\begin{lemma}[Adapted coordinates]\label{lem:B}
Let $n$ be the future unit normal to $S$, $(e_1=c',e_2,e_3)$ an $h$-orthonormal
frame along $c$ with $(n,e_1,e_2,e_3)$ positively oriented,
$\sigma(s,y)=\exp^h_{c(s)}(y_2e_2+y_3e_3)$, and
$\Phi(t,s,y)=\exp^g_{\sigma(s,y)}(t\,n)$. Write
$T_\rho=\{|t|<\rho,\ -\rho<s<L+\rho,\ |y|<\rho\}$. For some $\rho_0\le\delta_0$,
$\Phi$ is a diffeomorphism from $\overline{T_{\rho_0}}$ onto a closed
neighbourhood of $c([0,L])$, orientation- and time-orientation-preserving;
$\Phi(\{t=0\})\subseteq S$ while $\Phi(t,\cdot)\in\Ifut(S)$ for $t>0$ and
$\in\Ichr(S)$ for $t<0$; $\Phi^*g=-dt^2+h_t(s,y)$; and
$\|\Phi^*g-\eta\|\le K(|t|+|y|)$ on $\overline{T_{\rho_0}}$.
\end{lemma}

\begin{proof}
Gaussian normal coordinates to $S$ composed with Riemannian normal
coordinates of $S$ along the compact embedded arc; points $\Phi(t,\cdot)$ lie
on the normal timelike geodesic through a point of the acausal $S$. At
$(t,y)=(0,0)$ the metric is exactly $\eta$ because the frame is orthonormal on
$c$; the estimate is smoothness on a compact set.
\end{proof}

\begin{lemma}[Cone control]\label{lem:C}
Let $\eta_c=-c^2dt^2+ds^2+|dy|^2$, $\tfrac12<\mu<1<\lambda$, and let $k$ be a
symmetric bilinear form on $\mathbb R^4$. If
$\|k-\eta_1\|\le\varepsilon<\min\{(1-\mu^2)/(1+\mu^2),(\lambda^2-1)/(\lambda^2+1)\}$
(operator norm), then $k$ is Lorentzian with $\partial_t$ timelike, and for
$v=(1,u)$: $|u|\le\mu\Rightarrow k(v,v)<0$, $|u|\ge\lambda\Rightarrow
k(v,v)>0$. The bound is preserved by $\chi k+(1-\chi)\eta_1$, $0\le\chi\le1$.
\end{lemma}

\begin{proof}
$k(v,v)\le-1+|u|^2+\varepsilon(1+|u|^2)$ and $k(v,v)\ge(1-\varepsilon)|u|^2-
(1+\varepsilon)$. The spatial block is $\ge(1-\varepsilon)I>0$ and
$k(\partial_t,\partial_t)<0$, which forces signature $(-,+,+,+)$.
\end{proof}

\paragraph{Constants and the auxiliary spacetime.} In this order: fix
$\mu,\lambda,\varepsilon$ as in Lemma~\ref{lem:C}; Lemma~\ref{lem:B} gives
$\rho_0$ and $K$; choose $r>0$ with $3r<\rho_0$, $6Kr\le\varepsilon$ and
$8\lambda r<L$. Then $\overline{T_{2r}}\subset T_{3r}$; let $\chi$ be smooth,
$\chi=1$ on a neighbourhood of $\overline{T_{2r}}$,
$\operatorname{supp}\chi\subset T_{3r}$, and on
$\mathbb R^4$ put $g'=\chi\,\Phi^*g+(1-\chi)\eta$. On $T_{3r}$,
$\|\Phi^*g-\eta\|\le6Kr\le\varepsilon$, so Lemma~\ref{lem:C} applies to $g'$
everywhere and to $\Phi^*g$ on $T_{3r}$. Every $g'$-causal vector therefore
has $|d(s,y)/dt|\le\lambda$; $t$ increases strictly along future causal
curves; and an inextendible causal curve along which $t$ had a finite limit
would have convergent spatial coordinates and extend. Hence every level
$\{t=c\}$ is a Cauchy hypersurface and $M'=(\mathbb R^4,g')$ is globally
hyperbolic (MS Thm.~3.75), flat outside $T_{3r}$, and
$\psi=\Phi^{-1}:\Phi(T_{2r})\to T_{2r}$ is an isometry preserving orientation
and time orientation.

\begin{lemma}[The slab region is self-determined]\label{lem:D}
Let $B_0=\{s\in[-r/2,L+r/2],\ |y|\le r/2\}$, $\Sigma=\Phi(\{0\}\times B_0)\subset
S$ and $\Sigma'=\{0\}\times B_0$. Then:
\begin{enumerate}
\item[(i)] $D_M(\Sigma)\subseteq\Phi(T_{2r})$ and $D_{M'}(\Sigma')\subseteq T_{2r}$;
\item[(ii)] $\psi(D_M(\Sigma))=D_{M'}(\Sigma')$;
\item[(iii)] $N=\operatorname{int}D_M(\Sigma)$ is causally convex in $M$,
globally hyperbolic, connected, and contains $x,b$; likewise
$\psi(N)=\operatorname{int}D_{M'}(\Sigma')$ in $M'$.
\end{enumerate}
\end{lemma}

\begin{proof}
\emph{Escape paths.} Let $q\in\partial T_{2r}$ with $0<t(q)\le2r$. If the
projection $(s_q,y_q)$ is not in $B_0$, follow the vertical segment down to
$t=0$. Otherwise $q$ lies on the lid $t=2r$ with $|y_q|\le r/2$; move with
velocity $(dt,dy)=(-1,\tfrac\mu2\hat y)$ for coordinate time $2r$, where $\hat y$
is a unit vector ($y_q/|y_q|$ if $y_q\neq0$), ending at $|y|\in(r/2,3r)$
because $\mu r>r/2$ and $r/2+\mu r<3r$. Either path lies in $T_{3r}$, has
coordinate speed $<\mu$, so is timelike for $\Phi^*g$ and for $g'$
(Lemma~\ref{lem:C}), and ends at a point of $\{t=0\}$ outside $\Sigma'$, i.e.\ on
$S\setminus\Sigma$ in $M$ ($\Phi$ is injective). A past-directed causal curve
starting on a Cauchy hypersurface never returns to it. The time-reversed paths
serve points with $t(q)<0$.

(i) Let $p\in D^+_M(\Sigma)$ with $p\notin\Phi(T_{2r})$, and let $\gamma$ be a
past-inextendible causal curve from $p$. It meets $\Sigma$; let $\sigma_1$ be
its first point on $\Sigma$, and $q$ the last point of $\partial\Phi(T_{2r})$ on
the segment from $p$ to $\sigma_1$. After $q$ the segment lies in
$\Phi(T_{2r})$, where $t$ increases along future causal curves, so $t(q)>0$.
The segment $p\to q$, the escape path from $q$, and any past-inextendible
continuation form a past-inextendible causal curve from $p$ missing $\Sigma$, a
contradiction. Time-reverse for $D^-$. The same argument works in $M'$, where
$\{t=0\}$ is a Cauchy hypersurface.

(ii) Let $p\in D^+_M(\Sigma)$, so $p\in\Phi(T_{2r})$ by (i), and let
$\gamma'$ be a past-inextendible $g'$-causal curve from $\psi(p)$. If
$\psi(p)\in\Sigma'$ there is nothing to prove; otherwise $t(\psi(p))>0$. Along
$\gamma'$, $t\to-\infty$; let $\tau_0$ be the first parameter at which $t=0$ and
$\tau_{\rm exit}$ the first at which $\gamma'$ meets $\partial T_{2r}$. The chart
and the agreement $g'=\Phi^*g$ extend to a neighbourhood of
$\overline{T_{2r}}$, so segments of $\gamma'$ up to either time pull back by
$\Phi$. If $\tau_0\le\tau_{\rm exit}$ (equality included): were
$\gamma'(\tau_0)\notin\Sigma'$, the pulled-back segment ends on $S\setminus\Sigma$
and, continued past $S$, is a past-inextendible $M$-curve from $p$ missing
$\Sigma$ (points with $t>0$ are not on $S$). If $\tau_{\rm exit}<\tau_0$, then
$t>0$ at the exit point, and the pulled-back segment followed by the escape
path in $M$ again misses $\Sigma$. Both contradict $p\in D^+_M(\Sigma)$, so $\gamma'$
meets $\Sigma'$. The reverse inclusion is the same argument with $M$ and $M'$
exchanged, using $D_{M'}(\Sigma')\subseteq T_{2r}$: no global coordinate $t$ on
$M$ is needed, since a past-inextendible $M$-curve from a point of
$\Phi(T_{2r})$ either leaves $\Phi(T_{2r})$ or reaches the Cauchy surface $S$,
and the same two cases apply. $D^-$ is time-reversed.

(iii) For $q\in D^+(\Sigma)$, $\Jm(q)\cap\Jp(\Sigma)\subseteq D^+(\Sigma)$: a
past-inextendible curve from $r$ in this set, prefixed by a causal curve
$q\to r$, meets $\Sigma$; if only on the prefix, at $z$, then $\sigma\le r\le z$
for some $\sigma\in\Sigma$, and acausality gives $r=z\in\Sigma$. Hence $D(\Sigma)$
is causally convex: if $p,q\in D(\Sigma)$, $r$ lies on a causal curve $p\to q$,
and an inextendible causal curve $\lambda$ through $r$ misses $\Sigma$, then in
the mixed case $p\in D^-$, $q\in D^+$ the curves $(p\to r)\cdot\lambda^+$ and
$\lambda^-\cdot(r\to q)$ meet $\Sigma$ on $p\to r$ and on $r\to q$, and
acausality forces both meeting points to be $r$; the cases $p,q\in D^\pm$
follow from the inclusion just proved. For interior points choose $p^-\ll p$
and $q^+\gg q$ in $\operatorname{int}D$; then
$J(p,q)\subseteq\Ifut_{\mathrm s}(p^-)\cap\Ichr_{\mathrm s}(q^+)\subseteq
D(\Sigma)$, an open set (MS Props.~2.16, 2.31), so $N$ is causally convex, its causal diamonds are the
compact diamonds of $M$, and $N$ is globally hyperbolic (MS Def.~3.70). In
$M'$,
$D_{\eta_\lambda}(\Sigma')\subseteq D_{g'}(\Sigma')\subseteq D_{\eta_\mu}(\Sigma')$
(narrower cones give larger domains of dependence). The first inclusion puts
every relative-interior point of $\Sigma'$, in particular $x$ and $b$, in the
interior; the second puts the projection of every interior point with $t\neq0$
into $\operatorname{int}B_0$, so its vertical segment to $\{t=0\}$ lies in
$\Ichr(q^+)\cap\Ifut(z^-)\subseteq N'$ for suitable $q^+$, $z^-$, joining it
inside $N'$ to the connected set $\operatorname{int}\Sigma'$. Transfer by (ii).
\end{proof}

By Lemma~\ref{lem:D}, $N$ is an object of $\GH$ and the inclusion $N\hookrightarrow
M$ and $\psi|_N:N\to M'$ are morphisms. By (LC),
\begin{equation}\label{eq:transfer}
x\,B_M\,b\iff x\,B_N\,b\iff\psi(x)\,B_{M'}\,\psi(b).
\end{equation}

\begin{lemma}[Flat detour]\label{lem:F}
Let $P,Q$ lie in the flat region of $M'$, be $\eta_\lambda$-spacelike, and be
joined by a smooth embedded arc $\kappa$ contained in a hyperplane
$H=\{t=f(s)\}$ with $f$ affine, $|f'|<1/\lambda$, and at positive distance from
$\overline{T_{3r}}$. Then $P,Q$ are spacelike-separated in $M'$, and there is a
compact acausal slab $\Sigma''\subset H$, containing $P,Q$ in its relative
interior, such that $N''=\operatorname{int}D_{M'}(\Sigma'')$ lies in the flat
region, equals $\operatorname{int}D_\Mink(\Sigma'')$ under the identity, and is
an object of $\GH$ whose inclusions into $M'$ and $\Mink$ are morphisms. Hence
$P\,B_{M'}\,Q\iff P\,B_\Mink\,Q$.
\end{lemma}

\begin{proof}
A $g'$-causal curve between $P,Q$ would be $\eta_\lambda$-causal. Along
$g'$-causal curves $\tfrac{d}{d\tau}(t-f(s))\ge\dot t(1-\lambda|f'|)>0$ and
$t-f(s)\to\pm\infty$, so $H$ is a Cauchy hypersurface for $g'$ and for $\eta$,
crossed exactly once by every inextendible causal curve of either metric.
Extend $\kappa$ slightly past $P$ and $Q$ inside $H$, keeping it embedded and off $\overline{T_{3r}}$, and let $\Sigma''$ be its closed
$\delta$-neighbourhood in $H$. A boost along $s$ followed by a time translation, an isometry of $\eta$ preserving time orientation, makes $H$ horizontal: in the new coordinates $(\tilde t,z)$, $\eta=-d\tilde t^2+|dz|^2$ and
$H=\{\tilde t=0\}$. Put
\[
W^\pm=\{(\tilde t,z):\pm\tilde t\ge0,\ \overline B(z,|\tilde t|)\subseteq\Sigma''\},
\]
the two halves of $D_\Mink(\Sigma'')$. As $\delta\downarrow0$ the sets $\Sigma''$ decrease to the extended arc, which is compact with empty interior in $H$. So the largest radius $h(\delta)$ of a closed ball in $\Sigma''$ tends to $0$: a limit of balls of radius bounded below would lie in the arc. Take $\delta$ with $\delta+h(\delta)$ small compared with the $(\tilde t,z)$-distance from the extended arc to $\overline{T_{3r}}$. Every point of $W=W^+\cup W^-$ has $z\in\Sigma''$ and $|\tilde t|\le h(\delta)$, so $W$ is compact and has an open neighbourhood $U$ on which $g'=\eta$.
For $p$ with $\tilde t(p)\ge0$ let $C(p)=\Jm_\eta(p)\cap\{\tilde t\ge0\}$. If $p\in
W^+$ then $C(p)\subseteq W^+$, because the balls are nested.

$W^+\subseteq D^+_{M'}(\Sigma'')$: while a past-inextendible $g'$-causal curve from
$p\in W^+$ stays in $U$ and above $H$ it is $\eta$-causal, so it lies in the
compact set $C(p)\subset U$. Hence it reaches $H$ inside $C(p)\cap H\subseteq\Sigma''$.

$D^+_{M'}(\Sigma'')\subseteq W^+$: suppose $p\in D^+_{M'}(\Sigma'')\setminus W^+$, and let
$\gamma$ be a past-inextendible $g'$-causal curve from $p$. It crosses $H$ at a
point of $\Sigma''\subseteq W^+$. As $W^+$ is closed there is a last point of $\gamma$
before which it lies outside $W^+$, so we can take $q$ on $\gamma$ above $H$, with
$q\notin W^+$ and so close to $W^+$ that $C(q)\subset U$. Since $q\notin W^+$, a
straight $\eta$-causal segment in $C(q)$ runs from $q$ to a point of
$H\setminus\Sigma''$, and it is $g'$-causal. Then $\gamma$ up to $q$, this segment
and any past-inextendible continuation form a past-inextendible $g'$-causal curve
from $p$ that crosses $H$ only outside $\Sigma''$, a contradiction.

Time reversal handles $W^-$, so $D_{M'}(\Sigma'')=W=D_\Mink(\Sigma'')$. Its interior
$N''=\{|\tilde t|<\operatorname{dist}(z,H\setminus\Sigma'')\}$ lies in $U$, is connected
and contains $P,Q$. The argument of Lemma~\ref{lem:D}(iii) shows that it is causally
convex and globally hyperbolic in $M'$ and in $\Mink$.
\end{proof}

\begin{proof}[Proof of Theorem \ref{thm:main}]
(a) Suppose $B_\Mink\neq\U$; by Theorem~\ref{thm:flat}, $B_\Mink\subseteq\Jm$
relates no spacelike pair. Let $x\,B_M\,b$ with $x,b$ spacelike-separated, and
transfer by \eqref{eq:transfer}. Take $w=(-4r,0,0,0)$, $v=(4r,L,0,0)$ and the arc
$t(s)=-4r+(8r/L)s$, $y(s)=(10r\sin(\pi s/L),0)$, $0\le s\le L$: where $|t|\le3r$,
$s/L\in[1/8,7/8]$ and $|y|\ge10r\sin(\pi/8)>3r$, elsewhere $|t|>3r$, and the
slope $8r/L<1/\lambda$. Vertical segments are $\eta_\mu$-timelike, so
$w\ll'\psi(x)$ and $\psi(b)\ll'v$; (CB) gives $w\,B'\,\psi(x)$ and
$\psi(b)\,B'\,v$, and (T) gives $w\,B'\,v$. Lemma~\ref{lem:F} gives $w\,B_\Mink\,v$
with $w,v$ spacelike, a contradiction. If $x\,B_M\,b$ with
$x\in\Jp(b)\setminus\{b\}$, Lemma~\ref{lem:R} gives $e\ll x$ spacelike to $b$, and
$e\,B_M\,b$ by (CB) and (T), reducing to the previous case.

(b) Suppose $B_\Mink=\U$. For spacelike $x,b\in M$, transfer to $M'$ and use
$x'=(4r,0,0,0)$, $b'=(-4r,L,0,0)$ with the arc $t(s)=4r-(8r/L)s$, same $y(s)$:
$\psi(x)\,B'\,x'$ and $b'\,B'\,\psi(b)$ by (CB), and $x'\,B'\,b'$ by
Lemma~\ref{lem:F}, since $B_\Mink=\U$; so $x\,B_M\,b$. By Lemma~\ref{lem:E} and
(T), every pair is related.
\end{proof}


\section{Standpoint logic}\label{sec:logic}

Read $\Box A$ at $b$ as ``at every event $x\le b$, $A$'' on the frame $(M,R)$
with $bRx$ iff $x\le b$. If $A$ is a tenseless claim about what occurs at $x$,
this is the reflexive past modality of Diodorean tense logic, and the
standpoint reading of Fine and Lipman \citep{Fine2005,Lipman2023} adds no
formal content. $R$ is reflexive and transitive, so S4 is valid, and in
Minkowski spacetime the valid formulas are exactly those of S4.2
\citep{Goldblatt1980}. The axiom .2, $\Diamond\Box A\to\Box\Diamond A$, is valid at a
world for all valuations iff any two of its successors have a common
successor \citep{BdRV2001}. So .2 is valid at $b$ iff the subframe generated
by $b$ is directed, that is, iff any two events in $J^-(b)$ have a common
causal past.

In a model conformal to $(\eta_-,\eta_+)\times\mathbb R^3$, two events have a
common past iff $|X-Y|<\eta_x+\eta_y-2\eta_-$, so if $\eta_->-\infty$ the
directedness fails in the past of every event and .2 is valid at no event. This is the particle-horizon structure of Rindler
\citep{Rindler1956} in modal form, and it holds equally with or without an
inflationary era of finite conformal duration. The horizon problem of
cosmology \citep{Guth1981} is therefore not the failure of .2, which occurs
with and without inflation, but the failure of particular instances: pairs of points
on our last-scattering surface, at comoving separation up to
$2(\eta_0-\eta_{\rm ls})$, have no common past unless
$\eta_0-\eta_{\rm ls}<\eta_{\rm ls}-\eta_-$, which fails without inflation and
holds with enough of it. Two cautions. If persistence is imposed on suitable tenseless past-event
facts, rather than on centered claims such as ``this event is present now,''
the corresponding valuations are
constrained along the future direction, and the frame property that matters
may concern common futures, which these models have when $\eta_+=\infty$; the logic of persistent
valuations is a separate question. And Lewis and van der Schaaf study causal
modal logics of general spacetimes \citep{vdSchaaf2026}; we have not checked
how their results bear on the past-directed frames used here. Whether S4 is
exactly the logic of all past-directed globally hyperbolic frames we leave
open.


\section{Status of the results}\label{app:status}

``Reviewed'' means internal review with its record in the project repository, not refereeing. Every entry retains the
conditions stated in the text.

\begin{center}
\begin{tabularx}{\textwidth}{@{}lY@{}}
\toprule
Result & Status\\
\midrule
Theorem~\ref{thm:main} and corollaries & Reviewed, conditional on (LC), (T), (CB) and the off-shell premise. Lemma~\ref{lem:F} reproved 31 October with one model review. Expert read pending\\
Theorems~\ref{thm:class}, \ref{thm:weyl}, Proposition~\ref{prop:reg} & Reviewed. A planned independent null-boundary audit did not finish\\
Causal-set and energy-condition sections & Reviewed in the manuscript's own reviews. Separate audits were stopped unfinished\\
Theorem~\ref{thm:char}, models, Theorem~\ref{thm:diamonds} & Reviewed; proofs checked. Priority search bounded. Overlap with Geroch, Manchak and Earman noted\\
Theorem~\ref{thm:massive} & Reviewed research note; outline only here. Full proof: \texttt{research/horizons/mass-classification/RESULT.md} in the project repository. Existing lacuna results may imply it\\
Proposition~\ref{prop:mass} & Reviewed\\
Theorem~\ref{thm:map} & Assembled from the cited results; rows checked by one model review\\
Proposition~\ref{prop:norm} & Reviewed\\
Theorem~\ref{thm:same-past} & Finite-control audit accepted at bounded scope; regularity audited. Inherits external LOY inputs, including omitted proofs. Full proof: \texttt{research/lawful-completions/coordinator/} in the project repository. Priority comparison incomplete; much of the setting is in print\\
Lemma~\ref{lem:cones}, Theorem~\ref{thm:rigidity}, Proposition~\ref{prop:uncountable} & Bounded audit of a specialization. External uniqueness and common-extension theorems used at selected read depth\\
Section~\ref{sec:realms} & Proposition~\ref{prop:sp} elementary; horizon lemma and cylinder classification reviewed\\
Sections~\ref{sec:maximal}, \ref{sec:truth} & Coordinator arguments with bounded review; restricted counterinterpretation, not named theories\\
Section~\ref{sec:direction} & New 31 October. Proposition~\ref{prop:mirror} elementary. Albert, Price and Tooley read in part. Not yet red-teamed\\
Modal comb (Section~\ref{sec:prior}) & Internally checked counterexample at the abstract level. Later theory and priority separate\\
Parts of Section~\ref{sec:bounds} & One red-team, by the same model family as the author's assistant. Armstrong read in part. Fine, Broad, McCall read only through reports, Forrest through an abstract and one review, and Carlotto--Schoen and Kehle--Unger at summary depth\\
Section~\ref{sec:process} & Bounded review of the calculation\\
\bottomrule
\end{tabularx}
\end{center}

\begin{thebibliography}{99}
\bibitem[Aames(2022)]{Aames2022} J. Aames, Temporal becoming in a relativistic universe: causal diamonds and Gödel's philosophy of time, \emph{Eur.\ J.\ Philos.\ Sci.} 12 (2022) 44. doi:10.1007/s13194-022-00471-z.
\bibitem[Albert(2000)]{Albert2000} D. Z. Albert, \emph{Time and Chance}, Harvard University Press, 2000.
\bibitem[Andersson et al.(1998)]{AGH1998} L. Andersson, G. J. Galloway and R. Howard, The cosmological time function, \emph{Class.\ Quantum Grav.} 15 (1998) 309--322. doi:10.1088/0264-9381/15/2/006.
\bibitem[Armstrong(2004)]{Armstrong2004} D. M. Armstrong, \emph{Truth and Truthmakers}, Cambridge University Press, 2004.
\bibitem[Arthur(2006)]{Arthur2006} R. T. W. Arthur, Minkowski spacetime and the dimensions of the present, in D. Dieks (ed.), \emph{The Ontology of Spacetime}, Elsevier, 2006, ch.~7. doi:10.1016/s1871-1774(06)01007-2.
\bibitem[Arthur(2019)]{Arthur2019} R. T. W. Arthur, Relativity and the present, in \emph{The Reality of Time Flow}, Springer, 2019, ch.~6. doi:10.1007/978-3-030-15948-1\_6.
\bibitem[Belnap(1992)]{Belnap1992} N. Belnap, Branching space-time, \emph{Synthese} 92 (1992) 385--434. doi:10.1007/BF00414289.
\bibitem[Bernal and Sánchez(2005)]{BS2005} A. N. Bernal and M. Sánchez, Smoothness of time functions and the metric splitting of globally hyperbolic spacetimes, \emph{Commun.\ Math.\ Phys.} 257 (2005) 43--50.
\bibitem[Bernal and Sánchez(2006)]{BS2006} A. N. Bernal and M. Sánchez, Further results on the smoothability of Cauchy hypersurfaces and Cauchy time functions, arXiv:gr-qc/0512095.
\bibitem[Blackburn et al.(2001)]{BdRV2001} P. Blackburn, M. de Rijke and Y. Venema, \emph{Modal Logic}, Cambridge University Press, 2001.
\bibitem[Bolotin and Tanatarov(2013)]{BolotinTanatarov2013} Yu. L. Bolotin and I. V. Tanatarov, A thousand problems in cosmology: horizons, arXiv:1310.6329.
\bibitem[Brading(2015)]{Brading2015} K. Brading, Physically locating the present: a case of reading physics as a contribution to philosophy, \emph{Stud.\ Hist.\ Philos.\ Sci.\ A} 50 (2015) 13--19.
\bibitem[Broad(1923)]{Broad1923} C. D. Broad, \emph{Scientific Thought}, Kegan Paul, 1923.
\bibitem[Brunetti et al.(2003)]{BFV2003} R. Brunetti, K. Fredenhagen and R. Verch, The generally covariant locality principle, \emph{Commun.\ Math.\ Phys.} 237 (2003) 31--68. doi:10.1007/s00220-003-0815-7.
\bibitem[Callender(2000)]{Callender2000} C. Callender, Shedding light on time, \emph{Philos.\ Sci.} 67 (2000) S587--S599.
\bibitem[Cameron(2008)]{Cameron2008} R. P. Cameron, How to be a truthmaker maximalist, \emph{Noûs} 42 (2008) 410--421.
\bibitem[Carlotto and Schoen(2016)]{CarlottoSchoen2016} A. Carlotto and R. Schoen, Localizing solutions of the Einstein constraint equations, \emph{Invent.\ Math.} doi:10.1007/s00222-015-0642-4, arXiv:1407.4766.
\bibitem[Cheyne and Pigden(2006)]{CheynePigden2006} C. Cheyne and C. Pigden, Negative truths from positive facts, \emph{Australas.\ J. Philos.} 84 (2006). doi:10.1080/00048400600759092.
\bibitem[Choquet-Bruhat and Geroch(1969)]{CBG1969} Y. Choquet-Bruhat and R. Geroch, Global aspects of the Cauchy problem in general relativity, \emph{Commun.\ Math.\ Phys.} 14 (1969) 329--335. doi:10.1007/bf01645389.
\bibitem[Clifton and Hogarth(1995)]{CH1995} R. Clifton and M. Hogarth, The definability of objective becoming in Minkowski spacetime, \emph{Synthese} 103 (1995) 355--387. doi:10.1007/BF01089733.
\bibitem[Colyvan(1998)]{Colyvan1998} M. Colyvan, Can the Eleatic Principle be justified?, \emph{Can.\ J. Philos.} 28 (1998). doi:10.1080/00455091.1998.10715975.
\bibitem[Correia and Rosenkranz(2018)]{CR2018} F. Correia and S. Rosenkranz, \emph{Nothing To Come: A Defence of the Growing Block Theory of Time}, Synthese Library 395, Springer, 2018. doi:10.1007/978-3-319-78704-6.
\bibitem[Deng(2017)]{Deng2017} N. Deng, Making sense of the growing block view, \emph{Philosophia} (2017). doi:10.1007/s11406-017-9832-1.
\bibitem[Dieks(2006)]{Dieks2006} D. Dieks, Becoming, relativity and locality, in D. Dieks (ed.), \emph{The Ontology of Spacetime}, Elsevier, 2006, ch.~8. doi:10.1016/s1871-1774(06)01008-4.
\bibitem[Dorato(2006)]{Dorato2006} M. Dorato, The irrelevance of the presentist/eternalist debate for the ontology of Minkowski spacetime, in D. Dieks (ed.), \emph{The Ontology of Spacetime}, Elsevier, 2006.
\bibitem[Dorato(2011)]{Dorato2011} M. Dorato, The Alexandroff present and Minkowski spacetime: why it cannot do what it has been asked to do, in \emph{Explanation, Prediction, and Confirmation}, Springer, 2011. doi:10.1007/978-94-007-1180-8\_26.
\bibitem[Earman(1995)]{Earman1995} J. Earman, \emph{Bangs, Crunches, Whimpers, and Shrieks}, Oxford University Press, 1995.
\bibitem[Earman(2008)]{Earman2008} J. Earman, Reassessing the prospects for a growing block model of the universe, \emph{Int.\ Stud.\ Phil.\ Sci.} 22 (2008) 135--164. doi:10.1080/02698590802496680.
\bibitem[Fewster and Verch(2012)]{FV2012} C. J. Fewster and R. Verch, Dynamical locality and covariance: what makes a physical theory the same in all spacetimes?, \emph{Ann.\ Henri Poincaré} 13 (2012) 1613--1674, arXiv:1106.4785.
\bibitem[Fine(2005)]{Fine2005} K. Fine, Tense and reality, in \emph{Modality and Tense}, Oxford University Press, 2005, pp.~261--320.
\bibitem[Forrest(2006)]{Forrest2006} P. Forrest, General facts, physical necessity, and the metaphysics of time, \emph{Oxford Stud.\ Metaphys.} 2 (2006) 137--154. doi:10.1093/oso/9780199290581.003.0006.
\bibitem[Galloway et al.(2018)]{GLS2018} G. J. Galloway, E. Ling and J. Sbierski, Timelike completeness as an obstruction to $C^0$-extensions, \emph{Commun.\ Math.\ Phys.} (2018). doi:10.1007/s00220-017-3019-2.
\bibitem[Geroch(1977)]{Geroch1977} R. Geroch, Prediction in general relativity, in J. Earman, C. Glymour and J. Stachel (eds.), \emph{Foundations of Space-Time Theories}, Minnesota Studies in the Philosophy of Science 8, University of Minnesota Press, 1977, pp.~81--93.
\bibitem[Gödel(1949)]{Godel1949} K. Gödel, A remark about the relationship between relativity theory and idealistic philosophy, in P. A. Schilpp (ed.), \emph{Albert Einstein: Philosopher-Scientist}, 1949, pp.~557--562.
\bibitem[Goldblatt(1980)]{Goldblatt1980} R. Goldblatt, Diodorean modality in Minkowski spacetime, \emph{Studia Logica} 39 (1980) 219--236. doi:10.1007/BF00370321.
\bibitem[Grøn and Johannesen(2011)]{GronJohannesen2011} Ø. Grøn and S. Johannesen, FRW universe models in conformally flat spacetime coordinates. III: universe models with positive spatial curvature, \emph{Eur.\ Phys.\ J. Plus} 126 (2011) 30, arXiv:1103.4753.
\bibitem[Guth(1981)]{Guth1981} A. H. Guth, Inflationary universe: a possible solution to the horizon and flatness problems, \emph{Phys.\ Rev.\ D} 23 (1981) 347--356.
\bibitem[Hawking and Ellis(1973)]{HE1973} S. W. Hawking and G. F. R. Ellis, \emph{The Large Scale Structure of Space-Time}, Cambridge University Press, 1973. doi:10.1017/cbo9780511524646.
\bibitem[Impellitteri(2026a)]{ProjectPercolation} A. Impellitteri, Complex percolation on the principal and doubly-principal cylinder sigma-algebras, working manuscript, 2026.
\bibitem[Impellitteri(2026b)]{ProjectRole} A. Impellitteri, What work does existence do?, project note, 2026.
\bibitem[Kehle and Unger(2022)]{KehleUnger2022} C. Kehle and R. Unger, Gravitational collapse to extremal black holes and the third law of black hole thermodynamics, arXiv:2211.15742.
\bibitem[Kehrberger(2021)]{Kehrberger2021} L. M. A. Kehrberger, The case against smooth null infinity I: heuristics and counter-examples, arXiv:2105.08079.
\bibitem[Lewis and van der Schaaf(2026)]{vdSchaaf2026} M. Lewis and N. van der Schaaf, Some results on causal modalities in general spacetimes, arXiv:2601.14029v2.
\bibitem[Lipman(2023)]{Lipman2023} M. A. Lipman, Standpoints: a study of a metaphysical picture, \emph{J. Philos.} 120 (2023) 117--138.
\bibitem[Loss(2022)]{Loss2022} R. Loss, How the block grows, \emph{Am.\ Philos.\ Q.} 59 (2022) 377--389. doi:10.5406/21521123.59.4.04 (penultimate draft read).
\bibitem[Luk et al.(2016)]{LOY2016} J. Luk, S.-J. Oh and S. Yang, Solutions to the Einstein-scalar-field system in spherical symmetry with large bounded variation norms, arXiv:1605.03893v1.
\bibitem[MacBride(2024)]{MacBride2024} F. MacBride, Truthmakers, \emph{Stanford Encyclopedia of Philosophy}, 2024 revision.
\bibitem[Madarász et al.(2025)]{MSS2025} J. Madarász, M. Stannett and G. Székely, Definable coordinate geometries over fields, part 1: theory, arXiv:2507.10279.
\bibitem[Manchak(2008)]{Manchak2008} J. B. Manchak, Is prediction possible in general relativity?, \emph{Found.\ Phys.} 38 (2008) 317--321. doi:10.1007/s10701-008-9204-6.
\bibitem[Minguzzi and Sánchez(2008)]{MS2008} E. Minguzzi and M. Sánchez, The causal hierarchy of spacetimes, in \emph{Recent Developments in Pseudo-Riemannian Geometry}, EMS, 2008, pp.~299--358, arXiv:gr-qc/0609119.
\bibitem[Minguzzi(2015)]{Minguzzi2015} E. Minguzzi, The vacuum conservation theorem, \emph{Gen.\ Relativ.\ Gravit.} 47 (2015) 32, arXiv:1406.5915.
\bibitem[Molnar(2000)]{Molnar2000} G. Molnar, Truthmakers for negative truths, \emph{Australas.\ J. Philos.} 78 (2000) 72--86. doi:10.1080/00048400012349361.
\bibitem[Monton(2006)]{Monton2006} B. Monton, Presentism and quantum gravity, in D. Dieks (ed.), \emph{The Ontology of Spacetime}, Elsevier, 2006, ch.~14.
\bibitem[Müller(2006)]{Muller2006} T. Müller, On the problem of defining the present in special relativity: a challenge for tense logic, in \emph{Time and History}, De Gruyter, 2006. doi:10.1515/9783110333213.441.
\bibitem[Müller(2019)]{Muller2019} T. Müller, Defining a relativity-proof notion of the present via spatio-temporal indeterminism, \emph{Found.\ Phys.} (2019). doi:10.1007/s10701-019-00267-x.
\bibitem[Price(1996)]{Price1996} H. Price, \emph{Time's Arrow and Archimedes' Point}, Oxford University Press, 1996.
\bibitem[Putnam(1967)]{Putnam1967} H. Putnam, Time and physical geometry, \emph{J. Philos.} 64 (1967) 240--247. doi:10.2307/2024493.
\bibitem[Rakić(1997)]{Rakic1997} N. Rakić, Past, present, future, and special relativity, \emph{Brit.\ J. Philos.\ Sci.} 48 (1997) 257--280. doi:10.1093/bjps/48.2.257.
\bibitem[Rindler(1956)]{Rindler1956} W. Rindler, Visual horizons in world-models, \emph{Mon.\ Not.\ R. Astron.\ Soc.} 116 (1956) 662--677.
\bibitem[Ringström(n.d.)]{RingstromNote} H. Ringström, Existence of a maximal globally hyperbolic development, author note, https://people.kth.se/\textasciitilde hansr/mghd.pdf.
\bibitem[Ringström(2009)]{Ringstrom2009} H. Ringström, \emph{The Cauchy Problem in General Relativity}, European Mathematical Society, 2009.
\bibitem[Saunders(2002)]{Saunders2002} S. Saunders, How relativity contradicts presentism, \emph{R.\ Inst.\ Philos.\ Suppl.} 50 (2002) 277--292. doi:10.1017/S1358246100010602.
\bibitem[Savitt(2005)]{Savitt2005} S. Savitt, Time travel and becoming, \emph{The Monist} 88 (2005) 413--422. doi:10.5840/monist200588323.
\bibitem[Savitt(2006)]{Savitt2006} S. Savitt, Presentism and eternalism in perspective, in D. Dieks (ed.), \emph{The Ontology of Spacetime}, Elsevier, 2006.
\bibitem[Savitt(2009)]{Savitt2009} S. Savitt, The transient nows, in \emph{Quantum Reality, Relativistic Causality, and Closing the Epistemic Circle}, Springer, 2009. doi:10.1007/978-1-4020-9107-0\_18.
\bibitem[Savitt(2015)]{Savitt2015} S. Savitt, I $\heartsuit$ $\Diamond$s, \emph{Stud.\ Hist.\ Philos.\ Mod.\ Phys.} 50 (2015) 19--24. doi:10.1016/j.shpsb.2015.02.001.
\bibitem[Stein(1968)]{Stein1968} H. Stein, On Einstein--Minkowski space--time, \emph{J. Philos.} 65 (1968) 5--23. doi:10.2307/2024512.
\bibitem[Stein(1970)]{Stein1970} H. Stein, On the paradoxical time-structures of Gödel, \emph{Philos.\ Sci.} 37 (1970) 589--601. doi:10.1086/288328.
\bibitem[Stein(1991)]{Stein1991} H. Stein, On relativity theory and openness of the future, \emph{Philos.\ Sci.} 58 (1991) 147--167. doi:10.1086/289609.
\bibitem[Steward(2013)]{Steward2013} H. Steward, Processes, continuants, and individuals, \emph{Mind} 122 (2013) 781--812. doi:10.1093/mind/fzt080.
\bibitem[Stout(1997)]{Stout1997} R. Stout, Processes, \emph{Philosophy} 72 (1997) 19--27. doi:10.1017/S0031819100056631.
\bibitem[Tooley(1997)]{Tooley1997} M. Tooley, \emph{Time, Tense, and Causation}, Oxford University Press, 1997.
\bibitem[Tooley(2010)]{Tooley2010} M. Tooley, Time, truth, actuality, and causation: on the impossibility of divine foreknowledge, \emph{Eur.\ J. Philos.\ Relig.} 2 (2010). doi:10.24204/ejpr.v2i1.355.
\bibitem[Viglione(2026)]{Viglione2026} F. Viglione, Not-so-absolute cosmic simultaneity, \emph{Erkenntnis} (2026). doi:10.1007/s10670-026-01075-2.
\bibitem[Wald(1984)]{Wald1984} R. M. Wald, \emph{General Relativity}, University of Chicago Press, 1984.
\bibitem[Wawer(2021)]{Wawer2021} J. Wawer, Tensed metaphysics and non-local grounding of truth, \emph{Disputatio} (2021). doi:10.2478/disp-2021-0023.
\bibitem[Wüthrich and Callender(2017)]{WC2017} C. Wüthrich and C. Callender, What becomes of a causal set?, \emph{Brit.\ J. Philos.\ Sci.} 68 (2017) 907--925. doi:10.1093/bjps/axv040; arXiv:1502.00018.
\end{thebibliography}
\end{document}
